\documentclass{article}

% AMS packages and theorem style setup
\usepackage{amssymb, amsthm, amsmath, amsfonts, stmaryrd}

% Misc packages
\usepackage{faktor}
\usepackage{mathtools}
\usepackage{units}

\usepackage{hyperref}
\usepackage{cleveref}
\theoremstyle{definition}
    \newtheorem{example}{Example}[section]
    \newtheorem{nonexample}[example]{Non-Example}

    \newtheorem{definition}{Definition}[section]

\theoremstyle{remark}
    \newtheorem{remark}{Remark}[section]

\theoremstyle{plain}
    \newtheorem{proposition}{Proposition}[section]
    \newtheorem{theorem}{Theorem}[section]
    \newtheorem{lemma}{Lemma}[section]
    \newtheorem{corollary}{Corollary}[section]
    \newtheorem{exercise}{Exercise}[section]

\DeclareMathOperator{\Spec}{Spec}

% Typesetting shortcuts
\newcommand{\N}{\mathbb{N}}
\newcommand{\Z}{\mathbb{Z}}
\newcommand{\Q}{\mathbb{Q}}
\newcommand{\R}{\mathbb{R}}
\newcommand{\C}{\mathbb{C}}
\newcommand{\A}{\mathbb{A}}
\newcommand{\F}{\mathbb{F}}
\renewcommand{\P}{\mathbb{P}}
\newcommand{\spec}{\mathrm{spec}}
\renewcommand{\l}{\langle}
\renewcommand{\r}{\rangle}
\newcommand{\set}[1]{\{#1\}}

% Tikz
\usepackage{tikz-cd}
\usepackage{quiver}
\usepackage{colonequals}
\usepackage{microtype}
\title{The Geometry of Schemes}
\author{Students at ANU}

\begin{document}
\maketitle

\stepcounter{section}
\section*{Lecture 1: The animals before the taxonomy}
\begin{remark}
    An echidna is known by many different names in different languages: for example, \textit{nyingarn} (Noongar), \textit{bigibila} (Gamilaraay), \textit{tachyglossus} (scientific name), \textit{spiny anteater} (original English name). Some names are older than others; the Indigenous ones in particular have been around for tens of thousands of years. However, they all refer to the same underlying object, the animal itself, which has been around for much longer - likely tens of millions of years. Changing the name does not change the object. 
    
    We shall take this as an analogy for \textbf{schemes}. The objects we call schemes are, in many cases, already familiar to us. What is new is not the objects themselves, but rather the language and perspective through which we view them. \\
    
    Thus, we will postpone a formal definition for this lecture. Instead, we will examine a variety of examples that can be recognised as schemes. At first, you may already recognise them through different lenses: a set, a topological space, an algebraic variety. This is much like recognising an echidna as merely an animal - such an observation is correct, but it overlooks the richer structure that distinguishes the object. Our goal through these examples is to develop an intuition for schemes before we formally define them.
\end{remark}
\begin{example}[Affine space]
    In general we can work over any field, but for now we work over $\C$, the field of complex numbers.
    \begin{enumerate}
        \item Our first example of a scheme is the set of $n$-tuples of complex numbers. This is called \textbf{affine space}, and is denoted by
        \[
        \A^n=\{(a_1,\ldots,a_n)\mid a_i\in \C\}.
        \]
        \item We can also carve out subsets of affine space to make schemes:
        \[
        X\subset\A^2, \qquad X:=\{(x,y)\mid x^2+y^2-1=0\},
        \]
        \[
        Y\subset \A^3, \qquad Y=\{(x,y,z)\mid 27x^2-y+z=0,\; 51xyz=3,\; 3z^2-5y^2=9\}.
        \]
    \end{enumerate}
\end{example}
\begin{nonexample}
    The following is \textbf{not} a scheme:
    \[
    X\subset\A^2, \qquad X=\{y^2-e^x=0\} 
    \]
    Why not? Crucially, it uses an operation that cannot be generalised to any ring: the exponential map. When we say `generalised to any ring', we mean `uses the operations of a ring': addition, multiplication, and subtraction. This is what makes this set not an algebraic scheme - it uses something more than ring theory. 
\end{nonexample}
So far the examples given have all been schemes over $\mathbb{Z}$: the polynomials involved have all used integral coefficients, which make sense in any ring $R$ using the unique homomorphism $\mathbb{Z}\to R$. If for example, we defined the following:
\[
X\subset \A^2, \qquad X=\{(x,y)\mid ix^2-y^3=0\},
\]
this would be a scheme over $\C$, but not over $\mathbb{Z}$. That is, it makes sense in any ring that contains $\C$.
\begin{example}[Group scheme]
    We define
    \[
    \mathrm{GL}_n=\left\{\left(\begin{pmatrix}
        a_{11} &\cdots &a_{1n} \\
        \vdots & \ddots & \vdots \\
        a_{n1} &\cdots & a_{nn}
    \end{pmatrix},t\right) \Bigg|\;\; t\cdot \det(a_{ij})=1\right\}\subset \A^{n^2+1}
    \]
    This is what is called a \textbf{group scheme}. Indeed, we have maps
    \[
    \mathrm{GL}_n\times\mathrm{GL}_n\to\mathrm{GL}_n, \qquad ((a_{ij},t),(b_{ij},s))\mapsto \left(\sum_k a_{ik}b_{kj},st\right)
    \]
    \[
    \mathrm{GL}_n\to\mathrm{GL}_n, \quad (a_{ij},t)\mapsto (t\cdot \mathrm{Adj}(a_{ij}), \det(a_{ij}))
    \]
    Note that these maps are both `algebraic': they use only ring operations. Furthermore, there exists
    \[
    I=\left(\begin{pmatrix}
        1 &  & \\
        & \ddots & \\
        & & 1
    \end{pmatrix},1\right)\in\mathrm{GL}_n.
    \]
    Together, these satisfy the group axioms, leading to the term group scheme. Other examples of group schemes include $\mathrm{SL}_n$ and $\mathrm{SO}(n)$.
\end{example}
\begin{example}[Projective space]
    We now come to a less obvious example of a scheme, called \textbf{projective space}. For now we work over $\C$. We define
    \[
    \P^n=\{\text{One dimensional (vector) subspaces of }\C^{n+1}\}
    \]
    We can view this as the set of all nonzero elements of $\C^{n+1}$, modulo the equivalence relation that
    \[
    (a_0,\ldots,a_n)\sim\lambda(b_0,\ldots,b_n), \quad \lambda\in \C\setminus\{0\}.
    \]
    That is, two points are equivalent if they differ by a nonzero scalar. We denote the equivalence class of the point $(a_0,\ldots,a_n)$ by $[a_0:\cdots:a_n]$. The space $\P^n=\{[a_0:\cdots:a_n]\}$ turns out to be a scheme.
\end{example}
\begin{example}~\label{P1=x^2-yz}
    In projective space $\P^n$, we have subschemes. For example, in $\P^2$,
    \[
    C:=\{[X:Y:Z] \mid X^2-YZ=0\}
    \]
    is a subscheme. In fact, it will turn out to be `isomorphic' to $\P^1$ (whatever `isomorphic' means).
  \end{example}    
\begin{example}
  In $\P^5$, the following space is a subscheme:
    \[
    A=\{[a:b:c:k:l:m] \mid 4k^2=2b^2+2c^2-a^2, \; 4l^2=2c^2+2a^2-b^2, \; 4m^2=2a^2+2b^2-c^2\}
    \]
    This scheme is an example of something called a K3 surface.
\end{example}
\begin{nonexample}
    Consider the subset of $\P^2$, defined by
    \[
    \{[X:Y:Z] \mid X^2-Y^2Z^2=0\}.
    \]
    A quick inspection will show that this set is not actually well-defined: $(1,1,1)$ satisfies the equation, while $(2,2,2)$ doesn't, but they both represent the same equivalence class! This shows that not every polynomial defines a well-defined subset of projective space. It turns out that only homogeneous polynomials do, as scaling will multiply the entire expression by the same factor.
\end{nonexample}
\begin{example}
    Consider the subscheme of $\P^2$, given by
    \[
    C_X:=\{[X:Y:Z] \mid X\ne 0\}
    \]
    An \textbf{inequality} is just as valid for defining a scheme as an equality is, and this is a homogeneous equation, so we are in the clear. Since $X\ne0$, we have $[X:Y:Z]=[1:y:z]$, where $y=Y/X,z=Z/X$. Since $y,z$ can take any complex value, the set $C_X$ is in bijection with $\A^2$: we have a copy of affine space lying in $\P^2$. Looking at the complement,
    \[
    \P^2\setminus C_X=\{[0:Y:Z]\},
    \]
    which is easily seen to be a copy of $\P^1$.
    
    Taking analogously $C_Y$ and $C_Z$, we see that there are three copies of $\A^2$ lying in $\P^2$. Moreover, they cover $\P^2$, since $[0:0:0]$ is not a valid element. This displays a key property of schemes: a general scheme is one which has an open cover by \textbf{affine} schemes.
\end{example}
\begin{example}[Line bundles]
    We can define the following subset of $\P^n\times \A^{n+1}$:
    \[
    \Xi = \{([a_0:\cdots:a_n],b_0,\ldots,b_n) \mid (b_0,\ldots,b_n)\text{ lies on the line spanned by }(a_0,\ldots,a_n)\}
    \]
    Equivalent formulations of this condition include:
    \begin{itemize}
        \item $\mathrm{rank}\begin{pmatrix} a_0 &\cdots &a_n \\ b_0 & \cdots & b_n \end{pmatrix}\leq 1$
        \item $a_ib_j=a_jb_i \;\;\forall i,j$
    \end{itemize}
    In this sense, $\Xi$ is an algebraic scheme. We have a natural projection map $\pi$ to $\P^n$, and we are interested in its fibers. Indeed, the points that map to $\ell\in \P^n$ are exactly those $(b_0,\ldots,b_n)$ which lie on $\ell$. In this sense, $\Xi\to \P^n$ is called a \textbf{line bundle}, as its fibers are lines. This particular one is called the \textbf{tautological} or \textbf{universal} line bundle.
    
    Another example of a line bundle is the natural projection
    \[
    \P^n\times \A^1\to\P^n.
    \]
    Here the fibers are always $\A^1$, a line. This line bundle is the \textbf{trivial} line bundle. It turns out that it is not isomorphic to the universal line bundle!
\end{example}
\begin{remark}
  The notion of `isomorphic' has been brought up a few times now, so what do we mean when we say two schemes are isomorphic?
  Roughly speaking, a morphism between two schemes is a map defined only using algebraic operations (addition, multiplication, and subtraction).  An isomorphism is a morphism that admits an inverse morphism.
  \end{remark}
\begin{example}[Algebraic isomorphism]
  In Example~\ref{P1=x^2-yz} we asserted that $\P^1\cong\{X^2-YZ\}\subset\P^2$.
  Now we give the pair of inverse morphisms.

  In one direction, we have the map 
    \[
    \P^1\to \{X^2-YZ\}, \qquad [s:t]\mapsto [st:s^2:t^2].
    \]
    This map is a `morphism of schemes'. The inverse is given by
    \[
    \{X^2-YZ\}\to \P^1, \qquad [X:Y:Z]\mapsto \begin{cases}
        [Y:X] &\text{if }Y\ne0 \\
        [X:Z] &\text{if }Z\ne 0
    \end{cases} .
    \]
    Note that we never have $Y=Z=0$, since otherwise $X=0$ also. One can check that it is well-defined on the overlap $Y,Z\ne0$.
\end{example}
\begin{example}[Ring of algebraic functions]
    For this specific example, working over $\C$, let
    \[
    C=\{y^2-x^3-1=0\}\subset \A^2
    \]
    We have a natural projection to $\A^1$ given by $(x,y)\mapsto x$. We look at the ring of `algebraic functions' on $C$ - essentially this means polynomial functions. This turns out to be
    \[
    R=\C[x,y]/(y^2-x^3-1),
    \]
    since multiples of this polynomial vanish on $C$. We define a $\C$-algebra to be a ring $R$ with a homomorphism $\C\to R$. A $\C$-algebra homomorphism $R\to \C$ by definition does nothing to $\C$, and so is determined by the images $x\mapsto a,y\mapsto b$. These must satisfy $b^2=a^3+1$ - that is, $(a,b)\in C$. As a result, points of $C$ are in bijection with $\mathrm{Hom}_\C(R,\C)$, the ring of $\C$-algebra homomorphisms $R\to \C$.
    
    Similarly considering $\A^1$, its ring of algebraic functions is $\C[x]$, and points of $\A^1$ are in bijection with $\mathrm{Hom}_\C(\C[x],\C)$, as homomorphisms are determined by the image of $x$. 
    
    This bijection also translates to morphisms: given the map
    \[
    \C[x]\to \C[x,y]/(y^2-x^3-1)=R
    \]
    which sends $x\mapsto x$, we have a corresponding map
    \[
    C\cong \mathrm{Hom}(R,\C)\to \mathrm{Hom}(\C[x],\C)\cong \A^1,
    \]
    given by the projection $(a,b)\mapsto a$. That is, the morphism of schemes $C\to\A^1$ was reconstructed from the morphism of their corresponding rings, but in the opposite direction. 
\end{example}
\begin{remark}
    In time, to every ring $R$ we will associate a scheme called $\Spec R$. There will be a notion of `$k$-valued points' of $\Spec R$ (for $k$ a ring), and there will be a one-to-one correspondence between such $k$-valued points and $\mathrm{Hom}(R,k)$. When $k$ is an algebraically closed field, a $k$-valued point is called a \textbf{geometric point}.
\end{remark}
\begin{example}[Geometric points]
  Let $R=\Z$. We look for the $\C$-valued points of $\Spec\Z$, recalling that they are in one-to-one correspondence with $\mathrm{Hom}(\Z,\C)$. There is a unique homomorphism $\Z\to\C$, and hence only one $\C$-valued point. Similarly, there is only one $\overline{\F}_p$-valued point.
  
  We now consider $\Z[i]$. There is a unique homomorphism $\Z\to\Z[i]$, and by our above discussion we should expect a corresponding scheme morphism $\Spec\Z[i]\to\Spec\Z$, sending $k$-valued points to $k$-valued points. To find the $\C$-valued points of $\Spec\Z[i]$, we note that there are two homomorphisms $\Z[i]\to\C$, given by $i\mapsto i$ and $i\mapsto -i$. So, there are 2 distinct $\C$-valued points. For the $\overline{\F}_p$-valued points, for $p\ne2$ there are similarly $2$ homomorphisms, as there are two solutions to $x^2+1 = 0$ in \(\overline{\F}_p\). For $p=2$ there is only one homomorphism however (given by \(x = 1)\)
  
    As such, our corresponding morphism $\Spec\Z[i]\to\Spec\Z$ must send both $\C$-valued points to the single \(\C\)-valued point in $\Spec\Z$, and similarly for $\overline{\F}_p$ for $p\ne2$. For $p=2$, it sends the unique $\overline{\F}_2$-valued point to the unique one in $\Spec\Z$.
\end{example}


%%% Local Variables:
%%% mode: LaTeX
%%% TeX-master: "main"
%%% End:

% temporary newpage to make easier to write :)
\newpage


\stepcounter{section}
\section*{Lecture 2: details of sheaves and ringed spaces}

% TODO: 
% Remark about two ways to define manifolds, relate to schemes
% Define presheaf
% Examples presheaf
% Define sheaf
% Examples sheaf
% Remark sheaf with empty set
% Define ringed space
% Examples ringed spaces
% Define maximal ideal (???)
% Proposition about local ring / maximal ideal

\begin{remark}
    A manifold is a topological space equipped with extra data which forces it to ``locally resemble'' Euclidean space. Recall that there are two (equivalent) ways to define a manifold from a base topological space:
    \begin{enumerate}
        \item Using charts and atlases.
        \item By attaching a sheaf of functions.
    \end{enumerate}

    Usually one learns the chart definition and then derives and uses properties of sheaves while working with manifolds. We will work towards defining schemes in a similar manner, but primarily focus on the sheaf construction instead, and the analogues to charts will follow.
\end{remark}

We will need to define the machinery of sheaves to build towards schemes.


\begin{definition}[Presheaf]
    A \emph{presheaf} (of sets) $\mathcal{F}$ on a topological space $X$ consists of the following data:
    \begin{enumerate}
        \item For each open $U \subseteq X$, a set $\mathcal{F}(U)$.
        \item For each inclusion of open sets $V \subseteq U$, a \emph{restriction} function 
            \begin{align*}
            % Not sure what notation we're using so this for now , also should add map name to tikz diagram 
                \text{res}^U_V : \mathcal{F}(U) \to \mathcal{F}(V)
            \end{align*}
            which satisfies the following properties (contravariant functoriality):
            \begin{enumerate}
                \item For $V = U$, the restriction map for $V \subseteq U$ is the identity map $\text{res}_U^U = \text{id}_{\mathcal{F}(U)}$.
                \item For $W \subseteq V \subseteq U$, the following diagram commutes:
                    % https://q.uiver.app/#q=WzAsMyxbMCwwLCJcXG1hdGhjYWx7Rn0oVykiXSxbMiwwLCJcXG1hdGhjYWx7Rn0oVikiXSxbMiwyLCJcXG1hdGhjYWx7Rn0oVSkiXSxbMSwwXSxbMiwxXSxbMiwwXV0=
                    \[\begin{tikzcd}
                    	{\mathcal{F}(W)} && {\mathcal{F}(V)} \\
                    	\\
                    	&& {\mathcal{F}(U)}
                    	\arrow[from=1-3, to=1-1]
                    	\arrow[from=3-3, to=1-1]
                    	\arrow[from=3-3, to=1-3]
                    \end{tikzcd}\]
                
            \end{enumerate}
          \end{enumerate}
   A presheaf of groups/abelian groups/rings is a presheaf \(\mathcal{F}\) where \(\mathcal{F}(U)\) are groups/abelian groups/rings and where the restriction maps are homomorphisms.          
\end{definition}

We call the elements of \(\mathcal{F}(U)\) \emph{sections of \(\mathcal{F}\) on \(U\)} and the elements of \(\mathcal{F}(X)\) the \emph{global sections} of \(\mathcal{F}\).
For $f \in \mathcal{F}(U)$ and $V \subseteq U$, the restriction $\text{res}_V^U (f)$ is typically written  as $f|_V$, in line with restriction of functions.



\begin{example}\label{ex:presheaves}
    Let $X$ be a topological space. The following are examples of presheaves.

    \begin{enumerate}
    \item Take $X = \R^n$ with the usual Euclidean topology.
      Let \(\mathcal{F}\) be the presheaf whose sections on an open set \(U\) are the smooth functions \(f \colon U \to \mathbb{R}\), and the restriction maps are the usual restriction maps.
    \item 
       Let \(\mathcal{C}_X\) be the presheaf whose sections on an open set \(U\) are the continuous functions \(f \colon U \to \mathbb{R}\).
     \item This is a generalisation of the example above.  Replace the co-domain \(\mathbb{R}\) by any other topological space \(Y\).
     \item Suppose we have a topological space \(Y\) with a map \(\pi \colon Y \to X\).
       Let \(\operatorname{Sect}_{\pi}\) be the presheaf whose sections on \(U\) are continuous functions \(f \colon U \to Y\) such that \(\pi \circ f = \operatorname{id}_{U}\).
       This is called the sheaf of sections of \(\pi\).
              
        \item Fix a set $S$. For each open $U \subseteq X$, define $\mathcal{F}(U) = S$, and the restriction map to be $\text{id}_S$. This is the \emph{constant presheaf}.
        % \item Fix a topological space $Y$. The presheaf of maps into $Y$ is denoted
        %     \[
        %         \text{Maps}(-, Y) 
        %     \]
        % and let $U \subseteq X$ be open. Then the set of maps $f: U \to Y$, denoted $\text{Maps}(-,Y)$
        % \item 
    \end{enumerate}
    
\end{example}

% wanted to include footnote in case people forgot from analysis and for easy reference? probably not necessary.
\begin{definition}[Sheaf]

    Let \(X\) be a topological space and let  $\mathcal{F}$ be a presheaf on \(X\).
  We say that $\mathcal{F}$ is a \emph{sheaf} if for every open subset \(U \subset X\) and an open covering of \(U\) by open sets \(\{U_i\}_{i \in I}\), the following compatibility conditions are satisfied:
    \begin{enumerate}
        \item \textbf{Locality.} If $f, g \in \mathcal{F}(U)$ satisfy $f|_{U_i} = g|_{U_i}$ for all $i \in I$, then $f = g$. That is, ``sections are equal if they are locally equal''.
        \item \textbf{Gluing.} Given $f_i \in \mathcal{F}(U_i)$ such that $f_i|_{U_i \cap U_j} = f_j |_{U_i \cap U_j}$ for all $i, j \in I$, there exists a (unique) $f \in \mathcal{F}(U)$ such that $f|_{U_i} = f_i$ for all $i \in I$. 
    \end{enumerate}
\end{definition}

\begin{remark}
    The gluing condition only requires the existence of a section which restricts correctly to the open cover. However, uniqueness of this section is guaranteed by the locality condition.
\end{remark}

\begin{example} All of the presheaves in \Cref{ex:presheaves} are also examples of sheaves, except the last one (the constant presheaf).  If \(S\) contains more than one element, then the constant presheaf is not a sheaf (see the following proposition). 
\end{example}

However, not all presheaves are sheaves. The following is an important property of sheaves, which easily excludes certain presheaves from being sheaves.

\begin{proposition}[Sheaf behaviour on the empty set]
    If $\mathcal{F}$ is a sheaf, then $\mathcal{F}(\varnothing) = \{\cdot\}$ is a singleton set. 
\end{proposition}

\begin{proof}
    Let $f, g \in \mathcal{F}(\varnothing)$. The only open cover of $\varnothing$ is the empty covering. So $f = g$ (vacuously) on every part of the open cover. Then by the locality condition, $f = g$. So $\mathcal{F}(\varnothing)$ has exactly one element.
\end{proof}



\begin{nonexample}
    Fix $S$ to be a non-singleton set. The constant presheaf with $\mathcal{F}(U) = S$ is not a sheaf, because $\mathcal{F}(\varnothing) = S$ which is not a singleton.
\end{nonexample}

One might think we can fix this by manually setting $\mathcal{F}(\varnothing)$ to some singleton. But even this does not work in general, because it can break the gluing condition.

\begin{nonexample}
    Fix $X = \R$ and $S = \Z$. Take the constant presheaf, but modify it so that now $\mathcal{F}(U) = \{0\}$ (chosen arbitrarily to satisfy the singleton condition). In particular, take $U_1 = (0,1)$ and $U_2 = (2,3)$, which are disjoint and form an open cover for $U = U_1 \cup U_2$. Suppose $f_1 \in \mathcal{F}(U_1)$ is the constant function $1$, and $f_2 \in \mathcal{F}(U_2)$ is constant $-1$. These functions are compatible as they both restrict to the same element, $0$, on their intersection. But there is no function $f \in \mathcal{F}(U)$ which restricts correctly on both $U_1$ and $U_2$, because $1 \neq -1$.
\end{nonexample}

% It is also common for the topology of the underlying space to be too "non-algebraic" to admit 




% include something about locally constant being a valid sheaf?

Now we come to the defintion of a ringed space, which will be key moving forwards.

\begin{definition}[Ringed space]
    A \emph{ringed space} is a topological space $X$, equipped with a sheaf of rings $\mathcal{O}_X$.
\end{definition}

\begin{example}
  The topological space \(X = \mathbf{R}^n\) with \(\mathcal{O}_X = \mathcal{C}_X\) (the sheaf of continuous functions) is a ringed space.
  So is \(X = \mathbf{R}^n\) with \(\mathcal{O}_X = \mathcal{C}^{\infty}_X\) (the sheaf of smooth functions).
\end{example}

In future lectures we will be able to take a ring $R$, and construct from it a ringed space called $\text{spec}(R)$. Each $\text{spec}(R)$ will be an affine scheme, which we can glue together to form general schemes.



% An important property of ringed spaces is that they are entirely informed by the algebraic structure of a ring. 

\begin{nonexample}
    It is tempting to consider as an example, the space $X = \R$ with the presheaf as $\mathcal{F}(U) = \{\text{polynomial functions} : U \to \R \}$. But the gluing condition can break, for example if we take the zero polynomial on $U_1 = (-\infty, 0)$, and $x^2$ on $U_2 = (0, \infty)$. The polynomials automatically agree on the (empty) intersection, so should be compatible. But there is no global polynomial on the union $\R \setminus \{0\}$ which correctly restricts onto both $U_1$ and $U_2$.
\end{nonexample}

So what has gone wrong here, why is this not a sheaf? The answer is in the topology of the base space. The standard topology on $\R$ is the Euclidean topology, which relies on a number of non-algebraic operations such as square rooting, that may not be available or even compatible with rings. This motivates defining a topology which will be nicely compatible with all the algebraic properties we want to use. That will be another day. 

For now, let's build up some terminology to work with sheaves.

\begin{definition}[Stalk]
    Let $(X, \mathcal{F})$ be a sheaf. For a given point $x \in X$, we define the \emph{stalk} of $\mathcal{F}$ at $x$ to be 
    \begin{align*}
        \mathcal{F}_x & = \varinjlim_{U \ni x} \mathcal{F}(U) \\
        & = \bigsqcup_{U \ni x, f \in \mathcal{F}(U)} (U,f) \slash \sim
    \end{align*}
    where $(U,f) \sim (V, g)$ if there is a $W \ni x$, $W \subseteq U \cap V$, such that $f|_W = g|_W$. These equivalence classes are called \emph{germs}.
  \end{definition}

  If \(\mathcal{F}\) is a sheaf of groups, then the stalk of \(\mathcal{F}\) at a point \(x\) is a group.
  Similarly, for abelian groups, rings, etc.

The stalk captures the behaviour of the sheaf around a given point. 
\begin{example}
  Let $X = \R$ and $\mathcal{F} = \mathcal{C}_X$ the sheaf of continuous functions on $X$.
  The stalk $\mathcal{F}_0$ is a ring.
  It contains a copy of $\R$, corresponding to the constant functions, but it also contains other elements, represented by functions like $x, x^2, x^3, \sqrt{x}, e^x, \sin(x), |x|$ (all these represent distinct germs).

  The ring \(\mathcal{F}_0\) has no nilpotents, but is not an integral domain.
  It is a ``local ring'', meaning it has exactly one maximal ideal.
  The maximal ideal is given by the kernel of the evaluation map at zero:
    \[
        m = \ker (\text{eval}_0) = \{f \; | \; f(0) = 0\}.
    \]
    The evaluation map at zero is surjective onto $\R$ (consider the constant functions).
    So, the first isomorphism theorem gives $\mathcal{F}_0 \slash m \cong \R$.
    Since $\R$ is a field, this means $m$ is maximal.
    To see that \(m\) is the unique maximal ideal, we see that every element of \(\mathcal{F}_0\) that is not contained in \(m\) is a unit.
    Indeed, if \(f \in \mathcal{F}_0\) is such that \(f(0) \neq 0\), then \(f(x) \neq 0 \) for all \(x\) in a neighbourhood \(U\) of \(0\).
    Then the germ \((U, 1/f)\) is the multiplicative inverse of \(f\).
    Since the complement of \(m\) in \(\mathcal{F}_0\) consists of units, every non-unit ideal of \(\mathcal{F}_0\) must be contained in \(m\).
    So \(m\) is the only maximal ideal.
  \end{example}
  

% \begin{remark}
%     Arbitrary topologies (such as the standard Euclidean topology) may not behave according to the rigid properties of algebraic structures such as rings. 
% \end{remark}




% For now, let's become familiar with ringed spaces.

%%% Local Variables:
%%% mode: LaTeX
%%% TeX-master: "main"
%%% End:

\stepcounter{section}
\section*{Lecture 3: The spectrum of a ring}

\renewcommand\qedsymbol{$\mathbb{QED}$}

\begin{definition}
    A locally ringed space $(X, \mathcal{O}_X)$ is a topological space $X$, with a sheaf of rings $\mathcal{O}_X$ such that the stalks of $\mathcal{O}_X$ are local rings.
  \end{definition}

  For us, the most important example of a locally ringed space will be a scheme.
  The basic building blocks of schemes are affine schemes, which are locally ringed spaces \(\Spec R\) associated to a ring \(R\).
  The goal of this lecture is to define the topological space underlying \(\Spec R\).
  
\begin{definition}
    For a ring $R$, its spectrum $\Spec R$ is the set of all prime ideals of $R$. It has a topology called the Zariski topology, where all the closed sets are of the form $V(I)$, where $I$ is an ideal in $R$, and $V(I) = \{ p \in \Spec R : I \subseteq p \}$.
  \end{definition}
  \begin{remark}
    For \(f_1,\dots, f_n \in R\), we denote by \(V(f_1,\dots,f_n)\) the closed set \(V(I)\), where \(I\) is the ideal generated by \(f_1,\dots,f_n\).
    Note that \[V(f_1,\dots,f_n) = \{p \in \Spec R : f_i \in p \text{ for all } i\}.\]
  \end{remark}
\begin{theorem}
    The Zariski topology on the spectrum of a ring follows all the axioms of a topology, and the sets $D_f = \Spec R - V(f)$, for $f \in R$, form a basis for this topology.
\end{theorem}
\begin{proof}
    We first recognise that $V(0)$, the closed set associated with the zero ideal, is the entire space $\Spec R$, as every prime ideal contains the zero ideal. Similarly, $V(R)$ is the empty set, as no prime ideals can contain $R$, as then they'd be the entire space, which prime ideals are explicitly forbidden from being.

    All we have left to verify is that finite unions of $V(I)$ are closed and infinite intersections of $V(I)$ are closed.

    For finite unions, this is done by showing that $V(I) \cup V(J) = V(IJ)$. The forward inclusion is clear, as if $\mathfrak{p} \supset I$ or $\mathfrak{p} \supset J$, then $IJ$, which is contained in both $I$ and $J$, will surely be contained in $\mathfrak{p}$.
    For the reverse direction, $V(IJ) \subseteq V(I) \cup V(J)$ we need to show that if an arbitrary $\mathfrak{p}$ contains $IJ$, then it contains $I$ or $J$.
    Suppose \(\mathfrak{p} \not \supset I\), and let us prove that \(\mathfrak{p} \supset J\).
    Since \(\mathfrak{p} \not \supset I\), there exists \(i \in I\) such that \(i \not \in \mathfrak{p}\).
    Consider $ij$ for every $j \in J$. This product is in $IJ$, and so it is in $\mathfrak{p}$. Now as $\mathfrak{p}$ is prime, this means that either $i$ or $j$ is in $\mathfrak{p}$. As $i \notin \mathfrak{p}$, this means that for every $j \in J$ we have $j \in \mathfrak{p}$, so $J \subseteq \mathfrak{p}$. Thus, $V(IJ) \subseteq V(I) \cup V(J)$, and with the above, $V(IJ) = V(I) \cup V(J)$.
    
    For infinite intersections, consider $\bigcap_\alpha V(I_\alpha)$:
    \begin{align*}
        \bigcap_\alpha V(I_\alpha) &= \bigcap_\alpha \{ \mathfrak{p} \in \Spec \,R : I_\alpha \subseteq \mathfrak{p} \} \\
        &= \{ \mathfrak{p} \in \Spec \,R : \forall \alpha, I_\alpha \subseteq \mathfrak{p} \} \\
        &= \{ \mathfrak{p} \in \Spec \,R : \forall \alpha, \forall x \in  I_\alpha, x \in \mathfrak{p} \} \\
        &= \left\{ \mathfrak{p} \in \Spec \,R : \sum_i  x_i \in \mathfrak{p}, x_i \in I_i \right\} \\
        &= \left\{ \mathfrak{p} \in \Spec \,R : \sum_\alpha I_\alpha \subseteq \mathfrak{p} \right\} \\
        &= V\left(\sum_\alpha I_\alpha\right)
    \end{align*}
    The first line is by definition, the second moves the intersection inside the set definition. The third line unfolds the definition of a subset, and the fourth either adds terms of an ideal together, or takes only a single sum if going in the opposite direction. The fifth line is the definition of the sum of ideals, and the final is by definition again.
\end{proof}
To help understand the topologies of these spectra, we shall investigate a few examples. But first, a proposition that helps describes their topologies:
\begin{proposition}
    $R$ is a domain if and only if $(0) \in \Spec R$. If $(0) \in \Spec R$, then $V(0) = \Spec R$.
\end{proposition}
\begin{proof}
    $R$ being a domain is equivalent to $(0)$ being a prime ideal, and thus $(0) \in \Spec R$ by definition. For any ideal $I \subseteq R$, $(0) \subseteq I$ as $0 = 0\cdot i \in I$ for $i \in I$. Thus, as $\Spec R$ contains all the prime ideals of $R$, and all of them contain $(0)$, we get that $V(0) = \Spec R$.
  \end{proof}

  In all the examples, we describe all prime ideals \(\mathfrak{p}\) and the closed sets \(V(\mathfrak{p})\).
  Note that these are not all the closed sets! 
\begin{example}
    $\Spec (\mathbb{Z})$ consists of the elements $(p)$ where $p$ is a prime integer, and $(0)$. The closed set $V(p)$ for prime $p$ is $\{(p)\}$.  
\end{example}
\begin{proof}
  We know that all the non-zero prime ideals in the integers are those generated by prime integers, and $(0)$ as $\mathbb{Z}$ is a domain. As all the prime ideals of $\mathbb{Z}$ are maximal, the only prime ideal that can contain $(p)$ is $(p)$.
\end{proof}
\begin{example}
    $\Spec (\mathbb{Q}) = \{ (0)\}$.
\end{example}
\begin{proof}
    $\mathbb{Q}$ is a field, and thus the only non-zero ideal is $\mathbb{Q}$. Thus, the only prime ideal is $(0)$, as $\mathbb{Q}$ is not a prime ideal of $\mathbb{Q}$. There is only one possible choice of a topology on a one element set, so we are done.
\end{proof}
\begin{example}
    $\Spec (\mathbb{C}[x])$ contains the elements $(x - a)$ for all $a \in \mathbb{C}$, as well as the element $(0)$.  The set $V(0)$ is, again, $\Spec (\mathbb{C}[x])$ and $V(x-a)$ is the set $\{(x-a)\}$.
\end{example}
\begin{proof}
    We note that by Hilbert's Nullstellensatz all the maximal ideals of $\mathbb{C}[x]$ are of the form $(x-a)$ for $a \in \mathbb{C}$. It is quick to check that prime ideals must also be of this form. Pick a element of least degree from the a prime ideal $\mathfrak{p}$, if this degree isn't one, then we can find two other polynomials that multiply to it by the fundamental theorem of arithmetic. As $\mathfrak{p}$ is prime, one needs to be in $\mathfrak{p}$, so we have found a smaller degree polynomial in $\mathfrak{p}$, and thus there exists a degree $1$ polynomial in $\mathfrak{p}$.

    As there exists a degree $1$ polynomial in $\mathfrak{p}$, assuming monotonicity without loss of generality, say $x - a \in \mathfrak{p}$ for $a \in \mathbb{C}$. This means it contains the maximal ideal $(x-a)$, which means $\mathfrak{p} = (x-a)$ as $\mathbb{C}[x]$ is not a prime ideal. 

    For the topology, as all the prime ideals are maximal, the only prime ideal that will contain each is itself, so $V(x-a) = \{(x-a)\}$.). We also note that $\mathbb{C}[x]$ is a domain, so $(0)$ is in it's spectrum, and $V(0)$ is the entire spectrum.
\end{proof}
\begin{example}
  $\Spec (\mathbb{Z}/(24))$ contains the elements $(2)$ and $(3)$.
  We have $V(2) = \{(2)\}$ and $V(3) = \{(3)\}$.
\end{example}
\begin{proof}
    By the correspondence theorem, prime ideals of $\mathbb{Z}/(24)$ are in $1:1$ correspondence with prime ideals of $\mathbb{Z}$ containing $24$. The only two ideals that satisfy this condition are $(2)$ and $(3)$. We note that $(2)$ doesn't contain $(3)$, nor does $(3)$ contain $(2)$, they are distinct ideals in $\mathbb{Z}$ and $\mathbb{Z}/(24)$ too, so $V(2) = \{(2)\}$ and $V(3) = \{(3)\}$.
\end{proof}
\begin{example}
    $\Spec (\mathbb{C}[x] / (x^2)) = \{ (x) \}$.
\end{example}
\begin{proof}
    From the correspondence theorem, prime ideals of $\mathbb{C}[x]/(x^2)$ are in $1:1$ correspondence with prime ideals of $\mathbb{C}[x]$ containing $(x)$. As we have shown earlier, all the prime ideals of $\mathbb{C}[x]$ are of the form $(x-a)$ or they are $(0)$. The only one of these that contain $(x)$ is $(x)$ itself. So the only prime ideal of $\mathbb{C}[x]/(x^2)$ is the image of $(x)$, which is $(x)$. Again, there is only one choice of a topology on a one element set, so nothing more is needed.
\end{proof}
\begin{example}
    $\Spec (\mathbb{C}[x,y]/(y^2)) \simeq \Spec (\mathbb{C}[x])$. 
\end{example}
\begin{proof}
    We consider the element $y$ in $\mathbb{C}[x,y]/(y^2)$. It is nilpotent, and thus must be in any prime ideal in this ring. From the correspondence theorem, there is a $1:1$ bijection between prime ideals in $\mathbb{C}[x,y]/(y^2)$ containing $y$ and $\mathbb{C}[x,y]/(y^2)/(y) = \mathbb{C}[x]$. As all ideals in the domain contain this ideal, we thus get that $\Spec (\mathbb{C}[x,y]/(y^2)) \simeq \Spec (\mathbb{C}[x])$ as topological spaces. 
\end{proof}
\begin{example}
    $\Spec (\mathbb{C}[[t]])$ contains the elements $(t)$ and $(0)$. $V(t) = \{(t)\}$ and $V(0) = \Spec (\mathbb{C}[[t]])$.
\end{example}
\begin{proof}
    All elements in $\mathbb{C}[[t]]$ except for the powers of $t$ are invertible, so the only ideals that aren't the whole ring are the ideals $(0), (t), (t^2), \ldots$. From here, we can see that any power of $t$ greater than $1$ isn't prime, as $t\cdot t^{(n-1)}$ is a valid partition of it, and neither is in the ideal $(t^n)$. Thus, the only prime ideals of $\mathbb{C}[[t]]$ are $(0)$ and $(t)$. Now, we know $V(0)$ already, and it's clear to see that $V(t) = \{(t)\}$ as $(0) \subsetneq (t)$.
\end{proof}
\begin{example}
    $\Spec (\mathbb{Z}[\frac{1}{2}])$ contains all the elements $(p)$ where $p \neq 2$ and $(0)$. We also note that $V(0) = \Spec (\mathbb{Z})$ and $V(p) = \{p\}$ for a non-zero prime \(p\).
\end{example}
\begin{proof}
    We note that the injection that the integers has into $\mathbb{Z}[\frac{1}{2}]$ sends prime ideals of $\Spec (\mathbb{Z}[\frac{1}{2}])$ to prime ideals of $\Spec (\mathbb{Z})$. This is an observation that we prove a general fact about later. 
    
    We also note that it is a PID, as if we have $(m,n)$ an ideal, then we can multiply both by $2$ as many times as possible to ensure they are both in $\mathbb{Z}$, and then the integers being a PID gives us some $q$ such that $(m,n) = (q)$ through $am + bn = q$.

    Thus, all ideals in $\mathbb{Z}[\frac{1}{2}]$ have the form $(m)$ for some integer. Looking at their preimages in $\mathbb{Z}$, we see that the only possible prime ideals are those that map to the prime ideals in $\mathbb{Z}$, so $(0)$ and the odd primes. However, not the ideal $(2)$, as since $\frac{1}{2}$ exists in the domain, $(2) = \mathbb{Z}[\frac{1}{2}]$.

    Finally, we note that all these ideals are maximal again, as all the non-zero generating elements are co-prime, so $V(p) = \{p\}$ for non-zero $p$.
\end{proof}
\begin{example}
    $\Spec (\mathbb{Z}[\frac{1}{2},\frac{1}{3}])$ contains all the prime ideals elements $(p)$ where $p \neq 2$ or $p \neq 3$ and $(0)$. The sets $V(p)$ are the same as they would be for $\Spec (\mathbb{Z})$.
\end{example}
\begin{proof}
    Similarly to the above, we note that all the ideals have the form $(n)$ for some integral $n$, as we can simply multiply by $2$ and $3$ until the denominators are gone, and then collapse ideals generated by multiple elements into their collective gcd.

    Now, under the pre-image of the inclusion, the ideals that map to prime ideals are those generated by the prime integers other than $2$ and $3$, as well as $(0)$, so these are all the prime ideals. Similarly, they are all maximal, so $V(p) = \{(p)\}$ unless $p = 0$.
\end{proof}
\begin{definition}
    $\mathbb{Z}_{(2)}$ is the set of all fractions with odd denominator, or:
    \[ \mathbb{Z}_{(2)} = \left\{ x \in \mathbb{Q} : \exists p,q \in \mathbb{Z}, q \neq 0 \textit{ (mod 2)}, x = \frac{p}{q} \right\} \]
    This is called the localisation of $\mathbb{Z}$ away from $2$, and we will explore it more later.
\end{definition}

\begin{example}
    $\Spec (\mathbb{Z}_{(2)})$ contains the elements $(0)$ and $(2)$, with $V(2) = \{(2)\}$.
\end{example}
\begin{proof}
    Again, we get that all the ideals have form $(n)$ for integral $n$, and thus under the pre-image note that the only ideals map to prime ideals, and are not the entire ring are $(0)$ and $(2)$, as all the odd primes have their inverses included in the domain. We know what $V(0)$ is, and $V(2)$ is clearly just $\{(2)\}$ as $(0)$ doesn't include $2$.
\end{proof}

Now, we can make some observations from how we proved these statements.
Recall that the \emph{nilradical} of a ring, denoted by \(\sqrt{(0)}\) is the ideal containing all nilpotent elements.
\begin{theorem}
    The following observations are true:
    \begin{itemize}
        \item $\Spec (R) \simeq \Spec (R/\sqrt{(0)})$
        \item A ring morphism $\pi: R \longrightarrow S$ to induces a continuous morphism $\pi^{-1}: \Spec S \longrightarrow \Spec R$
        \item Closed points of \(\Spec R\) are the maximal ideals of \(R\).
    \end{itemize}
\end{theorem}
\begin{proof}
  Let \(a \in R\) be a nilpotent.
  Then \(a^n = 0\) for some \(n > 0\).
  So \(a^n \) lies in every prime ideal \(\mathfrak{p}\).
  But then \(a\) lies in every prime ideal \(\mathfrak{p}\).
  So, every prime ideal of \(R\) contains the nilradical.
  Thus, by the correspondence, there is a $1:1$ correspondence between prime ideals of $R$ and prime ideals of $R/\sqrt{(0)}$. This correspondence will respect the inclusions of ideals, and thus it is a homeomorphism of topological spaces, so $\Spec R \simeq \Spec (R/\sqrt{(0)})$

    Next, for $\pi : R \longrightarrow S$ and ring homomorphism, consider the pre-image of some prime ideal $\mathfrak{p} \in S$ ($\pi^{-1}(\mathfrak{p})$). This pre-image is definitely an ideal, and we can further note that it is the kernel of the induced map $\bar{\pi} : R \longrightarrow S / \mathfrak{p}$. Now, $S / \mathfrak{p}$ is a domain as $\mathfrak{p}$ is prime, so $R / \pi^{-1}(\mathfrak{p})$ is a domain, as it is isomorphic to the image of $\bar{\pi}$ by the first isomorphism theorem, and this image is in a domain, so it is a domain itself.

    As to showing that $\pi^{-1}$ is continuous, this was left as an exercise.

    For the closed points of the space:
    \[ \mathfrak{p} \textrm{ is closed} \iff V(\mathfrak{p}) = \{\mathfrak{p}\} \iff \mathfrak{p} \textrm{ is maximal} \]
\end{proof}

Now, all of these examples concern the spectra of rings. However, their maximal spectra, the space of a rings maximal ideals, a subset of it's spectrum, is very useful. They become the affine spaces that we were describing before.

\begin{example}
    $\mathrm{MaxSpec}(\mathbb{C}[x,y]/(y^2 - x^3 + x + 1))$ consists of all the ideals $(x-a, y-b)$ where $b^2 - a^3 + 1 + 1$, and $(a,b) \in \mathbb{C}$.
\end{example}
\begin{proof}
    Maximal ideals of $\mathbb{C}[x,y]/(y^2 - x^3 + x + 1)$ are the maximal ideals of $\mathbb{C}[x,y]$ that contain $(y^2 - x^3 + x + 1)$. From the Nullstellensatz, we know that all the maximal ideals of $\mathbb{C}[x,y]$ have the form $(x-a,y-b)$ for $(a,b) \in \mathbb{C}^2$.

    We note that quotienting by this ideal $(x-a, y-b)$ is effectively evaluating the polynomials in $\mathbb{C}[x,y]$ at $(a,b)$, and so we need that $y^2-x^3+x+1$ to evaluate to zero, so we require $b^2-a^3+a+1 = 0$.
\end{proof}

\renewcommand\qedsymbol{$\square$}

%%% Local Variables:
%%% mode: LaTeX
%%% TeX-master: "main"
%%% End:


\newpage

\stepcounter{section}
\section*{Lecture 4: Sheaf of Rings and Localisation}

In this lecture (number 4, originally given 6 August 2026), we define the Sheaf of Rings (and some elementary properties), and define and study the beginnings of Localisation.

Recall Grothenstieck: The topology carries the rings.

The Sheaf of Rings on the topological space of an affine scheme is called \emph{The Structure Sheaf}.
We recall from Lecture 3 that we need to associate a ring for every open set in $\Spec{R}$.
We showed that the empty set must be sent to a singleton, which for a ring means the zero ring.
We also recall that $\Spec{R}$ is the largest open set in $\Spec{R}$\footnote{
    Recall that the open sets are the complement of closed sets, so an open set is of the form $\Spec{R}\setminus{V(I)}$. Let $I=(1)$, then $V(I)$ is empty, so $\Spec{R}$ is the largest open subset.}, 
    so we can safely associate $\Spec{R}\rightsquigarrow R$.

This is supposed to be the analogue of smooth or holomorphic functions, so how do we consider the elements of $R$ as functions?

\begin{example}
    Let $R=\C[x,y]$, and note that these are polynomials over $x$ and $y$.
    We can consider the notion of evaluation of $f\in R$ by subsititution: the \emph{Value of $f$ at $p$} is $f(a,b)\in\C$ (where $f$ is a $\C$-valued function).
    This doesn't generalise well: how do you evaluate $67\in\Z$ (or any other member of $\Z$) in $\Spec\Z$? There is no polynomial for subsititution or evaluation.
\end{example}

Consider another approach that generalises well.
We recall that the underlying set (`points') of $\Spec{R}$ is the prime ideals of $\C[x,y]$ (consisting of $(0)$, ideals of irreducible polynomials in $\C[x,y]$ (see Hartshorne pg 7, Krull's \emph{Hauptidealsatz}, note that $\C[x,y]$ is a UFD), and the maximal ideals $(x-a,y-b)$ by the \emph{Nullstellensatz}).
For the maximal ideals, the map \((x-a,y-b) \mapsto (a,b)\) gives a bijection between the set of maximal ideals and \(\C^2\).
We now consider the image of $f\in\C[x,y]$ in $\faktor{\C[x,y]}{(x-a,y-b)}$.
As $(x-a,y-b)$ is a maximal ideal, this quotient gives us a field, and there's a canonical isomorphism of $\C$-algebras \(\C[x,y]/(x-a,y-b) \to \C\) given by $[f]\mapsto f(a,b)$.
In this way, we can think of the evaluation of \(f\) at \((a,b)\) as the image of \(f\) in the quotient ring \(\C[x,y]/(x-a,y-b)\). 
This way generalises to all rings and prime ideals, and we can think of elements of \(R\) as functions on \(\Spec R\).
\begin{example}
    Let $R=\C[x,y]$ and let $f\in\C[x,y]$.
    Then the evaluation at the maximal ideal $(x-a,y-b)$ corresponding to $(a,b)\in \C^2$ of $\Spec{R}$ is the image of $f$ in $\faktor{\C[x,y]}{(x-a,y-b)}$.
    As an example, consider $f=x-y$ at the points $(x-1,y-1)$ and $(x-1,y-2)$.
    Then at $(x-1,y-1)$, $f=(1-1)=[0]$ and at $(x-1,y-2)$ $f=(1-2)=[-1]$, which agrees with our intuition for evaluation.
\end{example}

\begin{example}
  Let $R=\Z$ and let $f=67$.
  Then the evaluation of \(f\) at a prime \(p\) is just the image of \(f\) in \(R/p\).
    For example, evaluating $67$ at $(2)$ gives $[1]$.
    The \emph{vanishing locus} of $V(67)$ is the set of all primes at which $67$ has value $[0]$, which is only $67$: $V(67)=(67)$.
    Similarly, $V(100)=\{(2),(5)\}$.
\end{example}

\begin{remark}
    $\Spec{\Z}$ is a really funky place to do geometry: $\Spec{\Z}$ describes a non-singular curve (Vakil).
  \end{remark}

  \begin{proposition}
    Let \(R\) be a ring.
    The closed points of \(\Spec R\) are the maximal ideals of \(R\).
  \end{proposition}
  \begin{proof}
    Recall that the prime ideals are the points of $\Spec{R}$ and that a maximal ideal is prime.
    Let \(\mathfrak{p}\) be a prime ideal such that \(\{\mathfrak{p}\}\) is closed.
    Then \(V(I) = \{\mathfrak{p}\}\) for some \(I\).
    As result, \(I \subset \mathfrak{p}\), and therefore \(V(I) \supset V(\mathfrak{p})\).
    So we get that \(\{\mathfrak p\} \supset V(\mathfrak{p})\).
    But clearly \(\{\mathfrak p\} \subset V(\mathfrak{p})\).
    So \(V(\mathfrak p) = \{\mathfrak p\}\).
    Note that \(V(\mathfrak p)\) is the set of all prime ideals containing \(\mathfrak p\).
    If \(V(\mathfrak p) = \{\mathfrak p\}\), then there is no prime ideal containing \(\mathfrak p\) other than \(\mathfrak p\) itself, which means that \(\mathfrak p\) is maximal.

    Conversely, if \(\mathfrak p\) is maximal, then \(V(\mathfrak p) = \{\mathfrak p\}\).
    This means that \(\{\mathfrak p\}\) is closed.
  \end{proof}
  
In traditional algebraic geometry, the rings we encounter are rings like
\begin{align*}
    &\faktor{\C[x,y]}{(y^2-x^3)}\\
    &\faktor{\C[x,y,z]}{(xy-z,zx-1)},
\end{align*}
or more generally rings like
\[ k[x_1,\dots,x_n]/I,\]
where \(k\) is an algebraically closed field and \(I \subset k[x_1,\dots,x_n]\) is an ideal.
\begin{proposition}
  The maximal ideals of \( k[x_1,\dots,x_n]/I\) are in bijection with the points \((a_1,\dots,a_n)\) of \(k^n\) such that for all \(f \in I\), we have \(f(a_1,\dots,a_n) = 0\).
  These are precisely the closed points of \(\Spec k[x_1,\dots,x_n]/I\).
\end{proposition}
For example, the maximal ideals of \(\mathbf{C}[x,y]/(y^2-x^3)\) are in bijection with the points \((a_1,a_2)\) of \(\C^2\) with \(a_2^2-a_1^3 = 0\).
\begin{proof}
  From the correspondence theorem between a ring and its quotient, we know that taking the pre-image gives a bijection between the maximal ideals of \(k[x_1,\dots,x_n]/I\) and the maximal ideals of \(k[x_1,\dots,x_n]\) that contain \(I\).
  From the Nullstellensatz, the maximal ideals of $k[x_1,\dots,x_n]$ are of the form \((x_1-a_1,\dots,x_n-a_n)\) for \((a_1,\dots,a_n) \in k^n\).
  Note that \(I \subset (x_1-a_1,\dots,x_n-a_n) \) if and only if for all \(f \in I\), we have \(f(a_1,\dots,a_n) = 0\).    
\end{proof}

We saw that we can think of elements of \(R\) as functions on \(\Spec R\).
However, it can happen that two distinct elements of \(R\) give the same function!
\begin{example}
    Let $R=\faktor{\Z}{24\Z}$. Take $f=12$.
    The points of $\Spec{R}$ are $(2)$ and $(3)$ because we have zero divisors (so $(0)$ is not prime).
    Note that $f(2)=f(3)=0$; so $f$ evaluated at any point of the scheme is zero
    But $f\neq 0$ (is distinct from the zero element of \(R\)).
\end{example}

This is a genuinely new feature, and can happen in a more traditional example.

\begin{example}
  Let $R=\faktor{\C[x]}{x^3}$ and $f=x$.
  The only point of \(\Spec R\) is the maximal ideal \((x)\).
  What is the value of $f$ at the point $(x)$?
  It is the image of \(f\) in \(R/(x)\), which is clearly $[0]$.
  So again a non-zero element of \(R\) gives the zero function on \(\Spec R\).
\end{example}

The previous examples show that evaluating $f$ at a point of $\Spec{R}$ is a lossy thing.

\subsection{Localisation}
At this point, we have assigned the ring $R$ to the largest open subset of $\Spec{R}$.
We need to assign rings to all open subsets.
    We need some new algebra to continue.

    Take a ring $R$, and a subset $S\subset R$ which is closed under multiplication.
    We will return to the examples of $S=\{1,f,f^2,f^3,\dots\}$ and $S$ the complement of a prime ideal (for example, $R=\Z$ and $S=(2)$, and $R_{(S)}$ is thus the odd numbers).
    Note that $S$ will always need to contain the unit; a short argument is that the product of zero elements must be $1$.
    If this is not satisfying, let us just make the convention that \(1 \in S\).
    
    We form a new ring: $S^{-1}R$ formed by equivalence classes of pairs \((r,s)\) where \(r \in R\) and \(s \in S\) modulo the equivalence relation $(r_1,s_1)\sim(r_2,s_2)$ if there exists $s\in S$ such that $s\cdot(r_1s_2-s_1r_2)=0$.

    If \(R\) is an integral domain and \(0 \not \in S\),  we simplify the equivalence relation as $(r_1,s_1)\sim(r_2,s_2)$ if
    \begin{equation}\label{eqn:olddef}
      (r_1s_2-s_1r_2)=0.
    \end{equation}
    In this case, it is fine to think of the elements of \(S^{-1}R\) as fractions \(r/s\).
    But the additional \(s\) in the definition of the equivalence relation is necessary in case \(R\) has zero-divisors.
    \begin{example}[{How the definition \eqref{eqn:olddef} breaks when $R$ not a domain}]
        We consider the equivalence relation defined by $(r_1,s_1)\sim(r_2,s_2)$ if $r_1s_2-s_1r_2=0$.
        Assume that $R$ is not an integral domain. Then there must exist zero divisors.
        We have that $(r_1,s_1)\sim(r_2,s_2)\implies r_1s_2=s_1r_2$, and $(r_2,s_2)\sim(r_3,s_3)\implies r_2s_3=s_2r_3$.
        Then we want to show that $(r_1,s_1)\sim(r_3,s_3)$.
        So we multiply through by $s_3$: $s_3r_1s_2-s_3s_1r_2=0$ and $s_1r_2s_3-s_1s_2r_3=0$.
        Note that we now have $s_3r_1s_2-s_1s_2r_3=0$, but if $s_2$ is a zero divisor, then $r_1s_3-s_1r_3$ may not be zero but may appear that way.
      \end{example}

      Even if \(R\) is not a domain, it is useful to think of the elements of \(S^{-1}R\) as fractions.
      We just have to be careful when we cross-multiply!

    We note that we have a homomorphism $R\to S^{-1}R$ that sends \(r\) to \((r,1)\) (recall that $S$ contains \(1\)).
    When does $a\neq0\in R\mapsto0$ in $S^{-1}R$?
    This happens if and only if $(a,1)\sim(0,1)$ and this occurs if and only if $\exists sa=0$, which implies that $S$ has got zero divisors.
    This shows that the map is not always injective.

    More generally, if $S$ contains no zero divisors, then the map is injective.
    This map is very rarely surjective.
    In fact, we will see that the localisations of \(R\) are the rings of functions on open subsets of \(\Spec R\).
    And we should not expect the map from global functions to functions on an open set to be surjective.  
    
    Finally, it can happen that the localisation is the zero ring, which occurs if and only if $0\in S$, which occurs if and only if $S$ contains a nilpotent.

    Under the natural map \(R \to S^{-1}R\), all the elements of $S$ are sent to units: the inverse of $(s,1)$ is $(1,s)$ (consider this written as a fraction).
    This is the smallest extension of $R$ such that all elements of $S$ become units.
    In other words, \(R \to S^{-1}R\) has the following universal property:
    \begin{equation*}
        \begin{tikzcd}
            R\ar[d]\ar[dr,"\varphi"]& \\
            S^{-1}R\ar[r,dashed,"\exists!"]&B
        \end{tikzcd}
    \end{equation*}
   If \(\phi \colon R \to B\) is any ring homomorphism such that images of elements of $S$ are invertible in $B$, then \(\phi\) factors uniquely through a map \(S^{-1}R \to B\).

    \subsection{Prime Ideals of Localisation}
        What are the prime ideals of $S^{-1}R$ and how does $\Spec{R}$ relate to $\Spec{S^{-1}R}$?
        \begin{equation*}
            \begin{tikzcd}
                R\ar[r,"i"]&S^{-1}R\\
                i^{-1}(\mathfrak{p})\ar[u, phantom, sloped, "\subset"]&\mathfrak{p}\ar[u, phantom, sloped, "\subset"]
            \end{tikzcd}
        \end{equation*}
        We first consider the map above, and note that $i^{-1}(\mathfrak{p})$ must be a prime ideal that does not contain an element of $S$, as elements of $S$ are sent to units in $S^{-1}R$ (by the correspondence theorem).
        In the other direction, let $\mathfrak{p}\in R$ be a prime ideal not containing any element of \(S\).
        Then $\l i(\mathfrak{p})\r=\left\{\frac{r}{s}\mid r\in\mathfrak{p}\right\} \subset S^{-1}R$ is a prime ideal.
        
        \begin{proposition}
          The maps \(\mathfrak p \mapsto i^{-1}(\mathfrak p)\) and \(\mathfrak p \mapsto \langle i(\mathfrak p)\rangle\) are mutually inverse bijections between the prime ideals of \(R\) disjoint from \(S\) and the prime ideals of \(S^{-1}R\).
        \end{proposition}
        \begin{proof}
          Let $\mathfrak{p}\in R$ be a prime ideal not containing any element of \(S\).
          Let us check that \(\langle i(\mathfrak p)\rangle\) is a prime ideal of \(S^{-1}R\).
          Suppose \(f_1/s_1 \cdot f_2/s_2 \in \langle i(\mathfrak p)\rangle\).
          Then \(f_1/s_1 \cdot f_2/s_2  = f/s\) for some \(f \in \mathfrak p\).
           By definition, $s^\prime sf_1f_2=s_1s_2fs^\prime$ (the equivalence relation), and is an equality in $R$.
            The RHS is a member of $\mathfrak{p}$, which gives us that the LHS must also be a member of $\mathfrak{p}$.
            But $s^\prime$ and \(s\) are not a members of $\mathfrak p$, which means that the other parts must be, so one of $f_1$ or $f_2$ must be in $\mathfrak{p}$
        \end{proof}

        What is $\Spec{R_f}$? (recall that $R_f=(1,f,f^2,\dots)^{-1}R$, the \emph{localisation of $R$ at $f$}).
        So we've got that $\Spec{R_f}=\Spec{R}\setminus{V(f)}\eqqcolon D_f$ (base of the topology), noting that $V(f)$ is the set of prime ideals where $f$ vanishes.
        So $\Spec{R_f}$ is the prime ideals of $R$ that don't contain $f$.

    \subsection{Building the Structure Sheaf}
        So we may now send $\Spec{R}\to R$ and $D_f\to R_f$ as the start of building the structure sheaf.
        If $D_f\subset D_g$, then we need to find a map $R_g\to R_f$.
        We need $(g)=(1)$ in $R_f$, otherwise we don't have a map.

        To show that \((g) = (1)\) in \(R_{f}\), we show that there is no prime ideal of \(R_f\) that contains \((g)\).
        Suppose there is.
        Take $\mathfrak{p}\subset R_f$ containing $g$ and look at the preimage $\tilde{\mathfrak{p}}\subset R$ of $\mathfrak{p}$.
        We note that $\tilde{\mathfrak{p}}$ contains $g$, but not $f$ as we're pulled back from a place where $f$ is invertible.  But this contradicts \(D_f \subset D_g\).

        We now have a pre-sheaf on the distinguished open sets, and this is the natural approach here.

        Recall the sheaf conditions:
        \begin{enumerate}
            \item If $f,g\in\mathcal{F}$ and $\{U_i\}$ is an open cover of some open $U$, then if $f\in\mathcal{F}(U_i)=g\in\mathcal{F}(U_i)$ for all $i$, then $f=g$
            \item Given $f_i\in\mathcal{F}(U_i)$ for all $i$ such that for all $i,j$ the image of $f_i\in\mathcal{F}(U_i\cap U_j)$ equals the image of $f_j\in\mathcal{F}(U_i\cap U_j)$ then there exists $f\in\mathcal{F}(U)$ that maps to all the $f_i$.
            \end{enumerate}

            It turns out that this procedure gives a ``sheaf on the base of the topology'' given by the open sets \(D_f\).
            By a general principle, this extends uniquely to a sheaf on the whole space.
            (This is explained in the next lecture).

        % If there is no prime ideal containig both $f$ and $g$, then $\Spec{R}=D_f\cup D_g$, then there are maps $R\to R_g$ and $R\to R_f$.
        % The first condition requires for every ideal that $f\neq 0$ or $g\neq 0$, as there is no ideal with $f=0$ and $g=0$.
        % Then $f\cup g=(1)$ and $(f,g)=(1)$, so we can find some $af+bg=1$.

        % Take $h_1$ and $h_2$ such that they are equal on both localisations.
        % So $(h_1-h_2)\cdot g^n=0$ and $(h_1-h_2)\cdot f^m=0$. Note that the differences are killed by powers of generators.
        % We want $h_1-h_2=0$ (but not as a function that evaluates to zero, actually zero).
        % For example, consider $n\cdot 2^m=0$ implies that $n=0$ (we don't necessarily have no nilpotents, so this is an important check).

        % Two sections that are equal locally must be equal globally.
        % We can check that the object that we have so far satisfies the sheaf axioms (Exercise.)

        % We now have the structure sheaf:
        % \begin{itemize}
        %     \item $\Spec{R}$ with this structure sheaf is our affine scheme
        %     \item Any open set can be covered by $D_f$s as they are a basis for the topology
        %     \item We need to check on double intersections
        %     \item This is effectively a book-keeping exercise
        % \end{itemize}

%%% Local Variables:
%%% mode: LaTeX
%%% TeX-master: "main"
%%% End:

\newpage

\stepcounter{section}
\section*{Lecture 5: Structure Sheaf of $\Spec R$}
Recall that for a ring $R$ we have defined an `incomplete' presheaf $\mathcal{F}$ on the base of the topology $\{D_f\}$, given by 
\[
\mathcal{F}(D_f) = R_f.
\]
We will show that this incomplete presheaf satisfies the sheaf axioms.
Let \(U\) be an open set in the base and \(\{U_i\}\) an open cover of \(U\) by open sets the base; that is, \(U = \bigcup U_i\).

We need to show that 
\begin{enumerate}
\item If $f,g\in \mathcal{F}(U)$ such that $f=g$ in $\mathcal{F}(U_i)$, for all $i$, then $f=g$ in $ \mathcal{F}(U)$.
\item If $f_i\in  \mathcal{F}(U_i)$ such that $f_i = f_j$ in $ \mathcal{F}(U_i\cap U_j)$, then there exists $f\in  \mathcal{F}(U)$ with $f=f_i$ in $ \mathcal{F}(U_i)$.
\end{enumerate}
Note that $U_i\cap U_j$ is not always in the base of the topology, in which case the condition $2.$ is replaced with

\begin{enumerate}
\item [2.] If $f_i\in  \mathcal{F}(U_i)$ such that $f_i = f_j$ in $ \mathcal{F}(V)$ for every \(V \subset U_i \cap U_j\) and \(V\)in the base, then there exists $f\in  \mathcal{F}(U)$ with $f=f_i$ in $ \mathcal{F}(U_i)$.
\end{enumerate}

However for this specific case we know that $D_f\cap D_g = D_{fg}$, so the base is closed under finite intersections. We leave the full verification of these axioms as an exercise, checking only the case $U = \text{spec }\!R$ with open cover $\text{spec }\!R = D_s\cup D_t$.
\begin{proof}[Proof of 1.]
We know that $\text{spec }\!R = D_s\cup D_t$ if and only if no prime ideal contains $s$ and $t$, so $\langle s,t\rangle = (1)$. By the assumption that $f=g$ in $R_s$ and $R_t$ we obtain
\[
s^n (f-g) = 0,\quad t^m(f-g) = 0
\]
for some $m,n\in\N$.

Since \(\langle s,t \rangle = (1)\), there is no prime ideal containing both \(s\) and \(t\).
Therefore, there is no prime ideal containing both \(s^{m}\) and \(t^n\).
In particular, there is no maximal ideal containing both \(s^{m}\) and \(t^{n}\).
So we must have \(\langle s^{m}, t^n \rangle = (1)\).
(This is a different proof of this fact that what we did in class.)

Take $A,B\in R$ such that $As^n + Bt^m = 1$, then
\[
0 = As^n (f-g) + Bt^m(f-g)= (As^n + Bt^m )(f-g) = f-g
\]
so $f=g$ in $R$ as required.
\end{proof}
\begin{proof}[Proof of 2.]
Let $f_1\in R_s$ and $f_2\in R_t$ be such that $f_1=f_2$ in $R_{st}$, we write
\[
f_1 = \frac{a}{s^{m}} = \frac{as^k}{s^{m+k}},\quad f_2 = \frac{b}{t^{n}} = \frac{bt^k}{t^{n+k}},
\]
where $k$ is chosen such that
\[
s^kt^k(t^na-s^mb) = 0.
\]
As before there exists $A,B\in R$ such that $As^{m+k} + Bt^{n+k} = 1$, so we let
\[
f = Aas^k + Bbt^k = \frac{Aas^k + Bbt^k }{As^{m+k} + Bt^{n+k}} \in R
\]
and claim that this is our desired $f$. To see that $f=f_1$ on $D_s$ observe
\begin{align*}
(Aas^k + Bbt^k)s^m - a(As^{m+k} + Bt^{n+k}) &= Aas^{m+k} + Bbt^ks^m - aAs^{m+k} - aBt^{n+k}\\
&= Bt^k(bs^m - at^{n}),
\end{align*}
with multiplying by $s^k$ yields $Bs^kt^k(bs^m - at^{n})=0$ as required. The argument on $D_t$ is much the same so we are done.
\end{proof}
We have shown that the axioms of a sheaf are satisfied on the base set, the following proposition will allow us to conclude that this defines a sheaf on the whole space as desired.
\begin{proposition}\label{prop:sheaves-on-a-base}
A presheaf on a base that satisfies the sheaf axioms on the base extends uniquely to a sheaf.
\end{proposition}
In fact the same is true for morphisms , that is, given sheaves on a base $\mathcal{F}, \mathcal{G}$ and a collection of maps 
\[
f(U): \mathcal{F}(U) \to \mathcal{G}(U)
\]
for every $U$ in our base, compatible with the restriction maps, then $f$ extends uniquely to a morphism of sheaves from the unique extension of \(\mathcal{F}\) to the unique extension of \(\mathcal{G}\).
The proposition could be stated as the following equivalence of categories 
\begin{equation}\label{eqn:base-cat}
\left\{
\begin{array}{c}
\text{sheaves on}\\
\text{a base}
\end{array}
\right\}
\;\simeq\;
\left\{
\begin{array}{c}
\text{sheaves on the}\\
\text{whole space}
\end{array}
\right\}.
\end{equation}
\begin{proof}[Proof of \Cref{prop:sheaves-on-a-base}]
  Suppose we are given a sheaf $\mathcal{F}$ on a base $\{U_i\}$.
  To extend it to a sheaf on the whole space, we need to specify \(\mathcal{F}(U)\) for an arbitrary open \(U\).
  The first candidate for the value of $\mathcal{F}(U)$ is to take some open cover $U = \bigcup U_i$ and define
    \[
    \mathcal{F}(U) = \left\{ \left(f_i \in \mathcal F(U_i)\right) \mid f_i = f_j\text{ on every basic }V\subset U_i \cap U_j \right\}.
  \]
  This works, but it is convenient because it involves choosing an open cover.
  For different choices of \(U\), we may have chosen different open covers, and then we need to do some work to compare them.
  To remove the choice, we can just take the biggest possible open cover. Thus we define
\[
\mathcal{F}(U) := \left\{ \left(f_i \in \mathcal{F}(U_i)\right)\,\forall U_i\subset U \mid \text{if } U_j \subset U_i \text{then } f_j \in \mathcal{F}(U_j)\text{ is the image of }f_i\in\mathcal{F}(U_i) \right\}
\]
Some might recognise this as the inverse limit $\varprojlim_{\substack{U_i \subset U}} \mathcal{F}(U_i)$. We leave it as an exercise to verify the sheaf axioms.
\end{proof}

\begin{example}
  The equivalence of categories in \eqref{eqn:base-cat} holds only for sheaves, and not for pre-sheaves.
  For example, consider the base of the Euclidean topology on \(\mathbf{R}^n\) consisting of open balls.
  Let \(\mathcal{F}\) be the pre-sheaf that sends every open ball to \(\mathbf{Z}\) (and the empty set to \(\{0\}\)).
  This is in fact a sheaf on the base, essentially because each open ball is connected.
  The unique extension of this sheaf is the sheaf of locally constant \(\mathbf{Z}\)-valued functions.
  However, this sheaf also has another extension to a presheaf, namely the pre-sheaf that sends every non-empty open set in \(\mathbf{R}^n\) to \(\mathbf{Z}\).
\end{example}

In our verification of the sheaf axioms for the incomplete structure sheaf, we used that \(D_s\) and \(D_t\) cover \(\Spec R\) if and only if \(\langle s,t\rangle = (1)\).
This observation leads to the following.
\begin{proposition}
  Suppose \(\Spec R = \cup_{i\in I} D_{s_i}\) for some potentially infinite set $I$.
  Then there is a finite subset \(J \in I\) such that \(\Spec R = \cup_{i\in J} D_{s_i}\).
  As a result, \(\Spec R\) is quasicompact (every open cover has a finite subcover).

\end{proposition}
\begin{proof}
  If \(\Spec R = \cup_{i\in I} D_{s_i}\), then no prime ideal contains all \(s_i\).
  Therefore, \(\langle s_i \rangle = (1)\).
  This means a finite linear combination of the $s_i$ is equal to $1$.
  Let \(J \subset I\) be the indices that appear in this linear combination.
   
  Since every open cover has a refinement by basic open sets, it follows that \(\Spec R\) is quasicompact.
\end{proof}

\begin{remark}
  The \emph{quasi} is used to signify that the space is not necessarily Hausdorff.
  In fact $\text{spec }\!R$ is almost never Hausdorff. Quasicompact spaces do not have the same advantages as compact ones, for example the continuous image of a qauasicompact space need not be closed. The better analogue for compactness is {\it properness}, which we will define later.
\end{remark}


From now on we will use $(X,\mathcal{O}_X)$ to denote the structure sheaf on $X =\Spec R$. The following examples will illustrate what the structure sheaf looks like in practice.
\begin{example}
    Let $R = \C[x,y]$, then we define $\A^2_\C:=\Spec R$ and consider the open set $U = \A^2 - (0,0) = \Spec R - (x,y)$, note $U$ is the complement of the closed point corresponding to the maximal ideal $(x,y)$. We wish to characterise $\mathcal{O}_{\A^2}(U)$. Note that $U$ consists of those prime ideals which do not contain either $(x)$ or $(y)$, thus
    \[
    U = D_x \cup D_y.
    \]
    We know that
    \[
    \mathcal{O}_{\A^2}(D_x) = \C[x,y]_x = \left\{ \frac{f(x,y)}{x^n} \mid f\in R, n\in\N\right\}
    \]
    and similarly
    \[
    \mathcal{O}_{\A^2}(D_y) = \left\{ \frac{g(x,y)}{y^m} \mid g\in R, m\in\N\right\}.
    \]
    Since $R$ is a domain we can regard all localisations as sitting inside the fraction field, in this case $\C(x,y)$. Hence
    \[
    \mathcal{O}_{\A^2}(U) = \left\{ \left(\frac{f(x,y)}{x^n}, \frac{g(x,y)}{y^m}\right) \text{ which are equal in }\C[x,y]_{xy} \right\}.
    \]
    However, because we are working in the fraction field the condition of equality is
    \[
    y^mf(x,y) - x^ng(x,y) = 0.
    \]
    So $f(x,y) = x^nF(x,y)$ and $g(x,y) = y^mG(x,y)$, and in fact $F=G$. Thus we may write
    \[
    \mathcal{O}_{\A^2}(U) = \left\{ (F, F) \right\}
    \]
    We then consider the map
    \[
    \mathcal{O}_{\A^2}(\A^2)\to \mathcal{O}_{\A^2}(U)  
    \]
    \[
    \C[x,y] \ni a \mapsto \{ (a,a) \}
    \]
    Which is an isomorphism of rings.
\end{example}
\begin{example}[Examples of structure sheaves]
We now give some more examples with familiar rings.
\begin{itemize}
    \item Let \(X = \Spec\C\). Since \(\C\) is a field,
    \[
        \Spec\C
        =
        \{(0)\}.
    \]
    Hence
    \[
        \mathcal{O}_X(X)
        =
        \C.
    \]
    \item Let
    \[
        X
        =
        \Spec\frac{\C[x]}{(x^2)}.
    \]
    The only prime ideal is \((x)\), since \(x\) is nilpotent. Thus
    \[
        X = \{(x)\},
    \]
    and
    \[
        \mathcal{O}_X(X)
        =
        \frac{\C[x]}{(x^2)}.
    \]
 \item Let \(X = \Spec\Q\). Then
    \[
        \Spec\Q
        =
        \{(0)\}.
    \]
    Hence
    \[
        \mathcal{O}_X(X)
        =
        \Q.
    \]
    \item Let
    \[
        X = \Spec(\mathbb{Z}/6\mathbb{Z}).
    \]
    The prime ideals are
    \[
        (2),\quad (3),
    \]
    so
    \[
        X = \{(2),(3)\}.
    \]
    We know
    \[
    \mathcal{O}_X(X) = \mathbb{Z}/6\mathbb{Z}
    \]
   Then
    \[
        D_2
        =
        \{(3)\},
    \]
    implies
    \[
        \mathcal{O}_X(D_2)
        =
        (\mathbb{Z}/6\mathbb{Z})_2
        \cong
        \mathbb{Z}/3\mathbb{Z}.
    \]
    Similarly,
    \[
        D_3 = \{(2)\} \implies \mathcal{O}_X(D_3) = (\mathbb{Z}/6\mathbb{Z})_3 \cong \mathbb{Z}/2\mathbb{Z}.
    \]
From which we recover $\Z/6\Z \cong \Z/2\Z\times \Z/3\Z$, as the Chinese remainder theorem states.
    \item Let
    \[
        X = \Spec\mathbb{C}[|t|].
    \]
    The prime ideals are
    \[
        (0),\quad(t),
    \]
    We know
    \[
     \mathcal{O}_X(X) = \C[|t|]
    \]
    and $D_0 = \emptyset\mapsto \{0\}$, thus we consider
    \[
    D_t = \{(0)\} \implies \mathcal{O}_X(D_t) = \C[|t|]_t = \C\left((t)\right)
    \]
\end{itemize}
\end{example}

Finally, let us compute the stalks of the structure sheaf.
Let $X = \Spec R$ and let \(\mathfrak{p}\in X\). Then the stalk at $\mathfrak{p}$ given by
    \[
        \mathcal{O}_{X,\mathfrak{p}}
        =
        \varinjlim_{\mathfrak{p}\in U}
        \mathcal{O}_X(U).
    \]
    Using the basis of open sets \(D_f\), this becomes
    \[
        \mathcal{O}_{X,\mathfrak{p}}
        =
        \varinjlim_{\substack{f\notin\mathfrak{p}}}
        \mathcal{O}_X(D_f)
        =
        \varinjlim_{\substack{f\notin\mathfrak{p}}}
        R_f
        \cong
        R_{\mathfrak{p}}.
    \]
    The notation \(R_{{\mathfrak{p}}}\) denotes the localisation of \(R\) in the multiplicative set that is the complement of \(\mathfrak p\).
\begin{example}
    In $\Spec \Z$ the stalk at $(5)$ is given by
    \[
    \left\{ \frac{a}{b} \mid 5\nmid b \right\}
    \] 
\end{example}
\begin{example}
    In $\Spec \C[x,y]$ the stalk at $(x,y)$ is given by
    \[
    \left\{ \frac{f(x,y)}{g(x,y)} \mid g(0,0)\neq 0 \right\}
    \] 
\end{example}
%%% Local Variables:
%%% mode: LaTeX
%%% TeX-master: "main"
%%% End:

\newpage

\stepcounter{section}
\section*{Lecture 6: Locally Ringed Spaces and their Morphisms}

\subsection{Two results about non-vanishing sections on a ringed space}

We begin with two lemmas, which will make essential use of the sheaf conditions. Recall for a locally ringed space $(X,\mathcal{O}_X)$ and $f \in \mathcal{O}_X(X)$,  the \textit{evaluation of $f$ at $x$} is the image of $f$ under the map $\mathcal{O}_X \to \nicefrac{\mathcal{O}_{X,x}}{\mathfrak{m}_x}$. We denote $f(x)$ as this evaluation.

\begin{theorem}
If $f(x) \ne 0$ for all $x \in X$, then $f \in \mathcal{O}_X(X)$ is a unit. 
\end{theorem}

\begin{proof}
  Because $f(x) \in \nicefrac{\mathcal{O}_{X,x}}{\mathfrak{m}_x}$ is non-zero, \(f \in \mathcal{O}_{X,x}\) does not lie in \(\mathfrak m_x\).
  Since \(\mathcal{O}_{X,x}\) is local, every element in the complement of the maximal ideal \(\mathfrak m_x\) is a unit.
  So there exists \(g_x \in \mathcal{O}_{X,x}\) such that \(f \cdot g_x = 1\) in \(\mathcal{O}_{X,x}\).
  This is to say that there is some function $g_x$ defined on an open neighbourhood $U_x$ of $x$ for which $f(x)g_x = 1$ on \(U_x\). Then we have an open cover $\{U_x\}_{x \in X}$ of $X$ and a multiplicative inverse \(g_x\) of \(f\) on each \(U_x\).

  We now use the sheaf axioms to prove that the \(g_x\) glue to give a global section \(g\).
  To do that, we consider the diagram
\begin{equation}
\begin{tikzcd}
& \mathcal{O}_X(X) \arrow[dl] \arrow[dr] &\\
g_x \in \mathcal{O}_X(U_x) \arrow[dr]&  & \mathcal{O}_X(U_y) \owns g_y\arrow[dl]\\
& \mathcal{O}_X(U_x \cap U_y)&
\end{tikzcd}
\end{equation}

Both $g_x|_{U_{x} \cap U_{y}}$ and $g_y|_{U_{x} \cap U_{y}}$ are inverses of $f|_{U_{x} \cap U_{y}} \in \mathcal{O}_X(U_x \cap U_y)$.
In any ring,  multiplicative inverses are unique, so $g_x|_{U_{x} \cap U_{y}} = g_y|_{U_{x} \cap U_{y}}$. 
By the sheaf axioms, the $g_x$s glue together to give a global $g$ which satisfies $fg = 1$ in $\mathcal{O}_X(X)$.
\end{proof}

\begin{theorem}
The set $Y = \{x \in X \, | \, f(x) \ne 0\} \subseteq X$ is open.
\end{theorem}

\begin{proof}
In the proof of Theorem 1.1, we found a local inverse $g_x$ of $f$ defined in some open neighbourhood $U_x$ of $x$ so long as $f(x) \ne 0$. But then $f$ is non-zero in this entire neighbourhood, so $U_x \subseteq Y$.
\end{proof}

\subsection{Morphisms of Locally Ringed Spaces}

Having defined locally ringed spaces, we now turn our attention to maps between them. As it turns out, by asking a few common-sense things of our morphisms, we will find that in the case when our locally ringed spaces are both spectra, the morphisms of locally ringed spaces are exactly in bijection with the ring homomorphisms on the underlying rings.

Let us consider what we will require from a map $(X,\mathcal{O}_X)$ to $(Y,\mathcal{O}_Y)$.
As both spaces have $X$ and $Y$ as underlying topological spaces, we will require a continuous map $\phi : X \to Y$. This does not uniquely define the sheaf data (e.g. from Assignment 1 in a ring $R$ containing nilpotent elements, evaluating $f$ at every point of $\Spec R$ does not uniquely determine $f$), so we will also require some way of comparing the sheaves on $X$ and $Y$. This is not possible a priori, as $\mathcal{O}_X$ and $\mathcal{O}_Y$ are sheaves of rings on distinct topological spaces. Nonetheless, we can compare the sheaves of rings by making use of the \textit{pushforward}. 

\begin{definition}
Given two topological spaces $X$ and $Y$, a continuous map {$\phi : X \to Y$} and a sheaf of rings $\mathcal{F}$ defined on $X$, the pushforward $\phi_\ast\mathcal{F}$ is a sheaf of rings on $Y$ defined by
\begin{equation}
(\phi_\ast\mathcal{F})(U) 	:= 	\mathcal{F}(\phi^{-1}U)
\end{equation}
for all open sets $U$ in $Y$.
\end{definition}

A map from \((X, \mathcal{O}_X)\) to \((Y, \mathcal{O}_Y)\) will consist of a continuous map \(\phi \colon X \to Y\) together with a morphism of sheaves of rings $\phi^\# : \mathcal{O}_Y \to \phi_\ast\mathcal{O}_X$.
We will require \(\phi^{\#}\) to satisfy an additional condition.
Roughly speaking, we will require \(\phi^{\#}\) to be compatible with evaluation maps.
We now make this precise.

For any continuous map \(\phi \colon X \to Y\) and any sheaf \(\mathcal{F}\) on \(X\), we have a natural map
\[ (\phi_{*} \mathcal{F})_{\phi(x)} \to \mathcal{F}_x.\]
This map sends an element of \(\phi_{*} \mathcal{F})_{\phi(x)}\) represented by a section \(s \in \phi_{*} \mathcal{F} (U)\) to the same section \(s \in \mathcal{F}(\phi^{-1}(U))\).
In particular, we have a map
\begin{equation}
(\phi_\ast \mathcal{O}_X)_{\phi(x)}  	\to 	\mathcal{O}_{X,x}. \label{FirstStalkMap}
\end{equation}
Composing this with the map $\phi^\#$ produces a ring homomorphism on stalks
\[\phi_x^\# : \mathcal{O}_{Y,\phi(x)} \to \mathcal{O}_{X,x}.\] 

The map \(\phi_x^{\#}\) allows us to compare the evaluation of \(f \in \mathcal{O}_{Y, \phi(x)}\) at \(\phi(x)\) with the evaluation of \(f^{\#} := \phi_x^{\#}(f)\) at \(x\).
This cannot be achieved in its entirety --- the evaluation of \(f\) lies in the field \(\mathcal{O}_{Y, \phi(x)}/ \mathfrak m_{\phi (x)}\) and the evaluation of \(f^{\#}\) lies in \(\mathcal{O}_{Y, \phi(x)}/ \mathfrak m_x\), and a priori there is no map between these quotient rings.
But it is possible to transfer the data that the respective functions evaluate to \(0\).
We require that \(f^{\#}(x) = 0\) if and only if \(f(\phi(x)) = 0\).
This condition motivates the following.
\begin{definition}
  Let \(R\) and \(S\) be local rings with maximal ideals \(\mathfrak{m}_R\) and \(\mathfrak{m}_{S}\).
A \emph{local ring homomorphism} is a ring homomorphism $f: R \to S$ for which $f^{-1}(\mathfrak{m}_S) = \mathfrak{m}_R$.
\end{definition}

Combining what we've done, we can now define a morphism of locally ringed spaces.

\begin{definition}
A morphism of locally ringed spaces $(X,\mathcal{O}_X) \to (Y,\mathcal{O}_Y)$, consists of the following:
\begin{enumerate}
\item a continuous map $\phi : X \to Y$, and
\item a morphism of sheaves $\phi^\# : \mathcal{O}_Y \to \phi_\ast\mathcal{O}_X$, where
\item for all $x \in X$, the induced map on stalks $\phi_x^\# : \mathcal{O}_{Y,\phi(x)} \to \mathcal{O}_{X,x}$ is a local homomorphism of local rings.
\end{enumerate}
\end{definition}

\begin{example}\label{Good Q to Z}
We will now construct a morphism of locally ringed spaces $(\Spec \Q, \mathcal{O}_{\Spec\Q}) \to (\Spec\Z, \mathcal{O}_{\Spec\Z})$ by providing the required maps.
\end{example}
Our continuous map $\phi : \Spec\Q \to \Spec\Z$ will be the map defined by sending the zero ideal of $\Q$ to the zero ideal of $\Z$. To determine our pushforward sheaf $\phi_\ast \mathcal{O}_{\Spec\Q}$, because every non-empty open set $U$ contains the zero ideal, we have
\begin{align*}
(\phi_\ast \mathcal{O}_{\Spec\Q})\left(U\right)
&= 	
\mathcal{O}_{\Spec\Q}\Big((0)\Big)\\
&= \Q.
\end{align*}
Then, writing the open sets as $U = D(p_1 \cdots p_n)$ we get
\begin{equation}
\mathcal{O}_{\Spec\Z}(U) = \Z[1/p_1 \cdots p_n].
\end{equation}

So let $\phi^\# : \mathcal{O}_{\Spec\Z} \to \phi_\ast\mathcal{O}_{\Spec\Q}$ be the inclusion map $\Z[1/p_1\cdots p_n] \xhookrightarrow{} \Q$ for each $U$.

Finally, let us check the last condition.
There is only one point to check, namely the zero ideal of \(\mathbb{Q}\).
At this point, we have
\[\phi_0^\# : \mathcal{O}_{\Spec\Z, (0)} \to \mathcal{O}_{\Spec\Q, (0)}\] and because
\begin{align*}
\mathcal{O}_{\Spec\Z, (0)} 	&= 	\Z_{(0)} = \Q, \text{ and }\\
\mathcal{O}_{\Spec\Q, (0)} 	&= 	\Q_{(0)} = \Q,
\end{align*}
the induced map is the identity map, which is a local homomorphism.

\begin{nonexample} \label{Bad Q to Z}
We will consider a continuous map from $\phi : \Spec\Q \to \Spec\Z$ and a morphism of sheaves for which the induced homomorphism on stalks is not a local homomorphism.

Consider the continuous map $\phi : \Spec\Q \to \Spec\Z$ defined by sending the zero ideal to $(p)$ for some (non-zero) prime $p \in \Z$. Again, let $U = D(1/p_1\dots p_n)$, then
\begin{equation}
(\phi_\ast\mathcal{O}_{\Spec\Q}) (U) 	= 	\begin{cases}
\{0\} 	&\text{if } p = p_i,\\
\Q &\text{otherwise},
\end{cases}
\end{equation}
and as before $O_{\Spec\Z}(\textnormal{Spec}\Z[1/p_1\cdots p_n]) = \Z[1/p_1\cdots p_n]$. We define $\phi^\#$ as either inclusion into $\Q$ or the unique map to the zero ring as appropriate, and we check this is a map of sheaf of rings by checking the diagram
\begin{equation}
\begin{tikzcd}
\Z\left[\frac{1}{p_1 \cdots p_n}\right] 	\arrow[r, hook] \arrow[d]
& \Q \arrow[d] \\
\Z\left[\frac{1}{p p_1 \cdots p_n}\right] \arrow[r, two heads] & \{0\}.
\end{tikzcd},
\end{equation}
This diagram obviously commutes because the homomorphism into the zero ring is unique. 

When we consider the induced map on stalks however, we find that
\begin{equation}
\mathcal{O}_{\Spec\Z,(p)} 
=
\Z_{(p)}
=
\left\{ \frac{m}{n} \, \Big| \, m,n \in \Z, p \!\not| \,m \right\}
\end{equation}
and again $\mathcal{O}_{\Spec\Q, (0)} = \Q$, so the induced map on stalks is inclusion $\Z_{(p)} \xhookrightarrow{} \Q$. But this is not a local homomorphism, as the preimage of the zero ideal in $\Q$ is not the ideal $(p)$ in $\Z_{(p)}$. In the context of our motivating remarks, we our induced map on stalks has destroyed the information that $p$ evaluates to 0 in $\nicefrac{\mathcal{O}_{\Spec\Z,(p)}}{\mathfrak{m}_p} = \nicefrac{\Z}{17}$.
\end{nonexample}

\begin{remark}
Example \ref{Good Q to Z} and Non-example \ref{Bad Q to Z} together show that there is a unique morphism of locally ringed spaces $(\Spec\Q, \mathcal{O}_{\Spec\Q}) \to (\Spec\Z, \mathcal{O}_{\Spec\Z})$, which in fact could be induced by the unique map $\Z \to \Q$. This is an example of a bigger theorem given below.
\end{remark}

\subsection{Maps of rings induce maps of spectra}
Let \(R\) and \(S\) be rings and \(f \colon R \to S\) a ring homomorphism.
We explain how \(f\) induces a map of locally ringed spaces
\[ \phi = \phi_{f} \colon \Spec S \to \Spec R.\] 
The map of topological spaces is defined by
\begin{equation}
 \phi \colon \mathfrak{p} \mapsto f^{-1}(\mathfrak{p})
\end{equation}
To define the map of sheaves
\[ \mathcal{O}_{\Spec R} \to \phi_{*} \mathcal{O}_{\Spec S},\]
it is enough to define it on basic open sets.
Consider the basic open \(D_a \subset \Spec R\).
It is easy to check that its pre-image in \(\Spec S\) is \(D_{f(a)}\).
Now, we have
\[ \mathcal{O}_{\Spec R} (D_{r}) = R_a \text{ and } \phi_{*} \mathcal{O}_{\Spec S} (D_a) =  \mathcal{O}_{\Spec S}(D_{f(a)}) = S_{f(a)}.\]
Since \(f \colon R \to S\) sends \(a \in R\) to a unit \(f(a) \in S_{f(a)}\), it induces a unique homomorphism
\[ R_a \to S_{f(a)}.\]
This gives the required map on sheaves of rings.
It can be checked that this is a local homomorphism on stalks.

\stepcounter{section}
\section*{Lecture 9: Closed Subschemes after a bit of Sheaf Theory}
Notes by Yuki Ishihara.

Given an abelian group \(G\) and a subgroup \(N \subset G\), we can construct the quotient group \(G/N\). This motivates us to wonder if we can construct the quotient sheaf of abelian groups.  So let $X$ be a topological space, $\mathcal{F}$ a sheaf on $X$ and $\mathcal{G} \subset \mathcal{F}$ a subsheaf of abelian groups. One approach to construct $\mathcal{F}/\mathcal{G}$ is to send the open set $U \to \mathcal{F}(U)/\mathcal{G}(U) =: \mathcal{F}/\mathcal{G}(U)$.  This is a good candidate for the quotient and is a good presheaf however is not a sheaf in general. The following example illustrates how the gluing axiom can fail.
\begin{example}
    $X = \mathbb{P}_\C^1, \mathcal{F} = \mathcal{O}_X$ and $G \subset \mathcal{O}_X$ subsheaf consisting of sections that evaluate to 0 at 0 and $\infty$. Or:
    $$\mathcal{G}(U) = \left\{f \in \mathcal{O}_X : f(0) = 0 \text{ if } 0 \in U \text{ and } f(\infty) = 0 \text{ if } \infty \in U\right\}$$
    We consider the points $0 = [0 : 1] \in \mathbb{P}_\C^1$ and $\infty = [1 : 0] \in \mathbb{P}_\C^1$ and the open set $U = \mathbb{P}^1_\C$. Let $U_1 = U\backslash \{\infty\}$ and $U_2 = U \backslash \{0\}$ an open covering of $U$. \smallskip \\
    On $U$ we have $\mathcal{O}_X(U) = \C$ the constant functions and $\mathcal{G}(U) = \{0\}$ ie. the only constant map that evaluates to 0 at 0 and $\infty$ is the zero map. So $\mathcal{O}_X/\mathcal{G}(U) = \C$ \smallskip \\
    On $U_1 = \Spec\C[x]$ we have $\mathcal{O}_X(U_1) = \C[x]$ and $\mathcal{G}(U_1) = x\C[x]$. So $\mathcal{O}_X/\mathcal{G}(U_1) = \C$. \smallskip \\
    On $U_2 = \Spec\C[x^{-1}]$ we have $\mathcal{O}_X(U_2) = \C[x^{-1}]$ and $\mathcal{G}(U_2) = x^{-1}\C[x^{-1}]$. So $\mathcal{O}_X/\mathcal{G}(U_2) = \C$ \smallskip \\
    On the intersection $U_1 \cap U_2 = \mathbb{P}^1_\C$ we have $\mathcal{O}_X(U_1 \cap U_2) = \C[x, x^{-1}]$ and $\mathcal{G}(U_1 \cap U_2) = \C[x, x^{-1}]$. So $\mathcal{O}_X/\mathcal{G}(U_1 \cap U_2) = 0$. \smallskip \\
    Now take $1 \in \mathcal{O}_X/\mathcal{G}(U_1)$ and $2 \in \mathcal{O}_X/\mathcal{G}(U_2)$. These both restrict to 0 on the intersection trivially. The gluing axiom requires us to find a global section represented by a complex number $f = c$ that restricts to $1$ and $2$ on $U_1$ and $U_2$ respectively. However the restrictions $f|_{U_1} = c$ and $f|_{U_2} = c$ are just the inclusion map and since $1 \neq 2$ there cannot exist such an $f$.
\end{example}
To overcome this issue we use a tool known as sheafification. This follows from the following theorem.
\begin{theorem} (Sheafification) Given a presheaf $\mathcal{P}$ on $X$ there exists a sheaf $\tilde{\mathcal{P}}$ (called the sheafification of $\mathcal{P}$) on $X$ and a map $p : \mathcal{P} \to \tilde{\mathcal{P}}$ such that:
\begin{enumerate}
    \item For any sheaf $\mathcal{Q}$ on $X$ and a map $q : \mathcal{P} \to \mathcal{Q}$ there is a unique map $\tilde{q} : \tilde{\mathcal{P}} \to \mathcal{Q}$ such that the following diagram commutes
        \[\begin{tikzcd}
	{\mathcal{P}} && {\tilde{\mathcal{P}}} \\
	\\
	&& {\mathcal{Q}}
	\arrow["p", from=1-1, to=1-3]
	\arrow["q"', from=1-1, to=3-3]
	\arrow["{\tilde{q}}", from=1-3, to=3-3]
\end{tikzcd}\]
\item For any $x \in X$, the map \(\mathcal{P} \to \tilde{\mathcal{P}}\) induces an isomorphism on stalks $\mathcal{P}_x \cong \tilde{\mathcal{P}}_x$.
\end{enumerate}
\end{theorem}
\newpage 
Considering this we can define:
$$\mathcal{F}/\mathcal{G} := \text{Sheafification of } (U \mapsto \mathcal{F}/\mathcal{G}(U))$$
as the quotient sheaf of abelian groups. It comes along with the following universal property:
\[\begin{tikzcd}
	{\mathcal{G}} && {\mathcal{F}/\mathcal{G}} \\
	& {\mathcal{F}} \\
	& {\mathcal{Q}}
	\arrow["0", from=1-1, to=1-3]
	\arrow["i", from=1-1, to=2-2]
	\arrow["0"', from=1-1, to=3-2]
	\arrow["{\exists!}", from=1-3, to=3-2]
	\arrow["\pi", from=2-2, to=1-3]
	\arrow["q"', from=2-2, to=3-2]
\end{tikzcd}\]
where $i$ is the inclusion map and $\pi$ is the canonical quotient map.
That is, given a sheaf \(\mathcal{Q}\) and a map \(\mathcal{F} \to \mathcal{Q}\) such that the composite \(\mathcal{G} \to \mathcal{Q}\) is zero, there exists a unique map \(\mathcal{F}/ \mathcal{G} \to \mathcal{Q}\) as in the diagram above.

Motivated by group theory we now aim to reproduce the first isomorphism theorem in terms of sheaves. In order to do this we first need to know the following definitions.
\begin{definition} The following are definitions on maps $\mathcal{F} \to \mathcal{Q}$ between sheaves of abelian groups. 
\begin{itemize}
    \item (Surjectivity) The map $\mathcal{F} \to \mathcal{Q}$ is surjective is it is surjective on all stalks. 
    \item (Kernel) The kernel sheaf $\mathcal{K}$ of the map $\mathcal{F} \to \mathcal{Q}$ is the sheaf $\mathcal{K}(U) = \ker(\phi_U)$ where $\phi_U : \mathcal{F}(U) \to \mathcal{Q}(U)$ is the induced map. 
\end{itemize}
\end{definition}
\begin{example}
    The quotient map $\mathcal{F} \to \mathcal{F}/\mathcal{G}$ is surjective sending the stalk $\mathcal{F}_x \mapsto \left[\mathcal{F}_x\right]$. 
\end{example}
We then get the first isomorphism equivalent for sheaves. 
\begin{corollary} (First Isomorphism Theorem)
    If $\mathcal{F} \to \mathcal{Q}$ surjective and $\mathcal{K}$ the kernel sheaf then $\mathcal{F}/\mathcal{K} \to \mathcal{Q}$ is an isomorphism of sheaves. 
\end{corollary}
\begin{exercise}
    Any map on sheaves that is an isomorphism on the stalks is an isomorphism
\end{exercise}
Caution: $\mathcal{F} \to \mathcal{Q}$ surjection of sheaves does not mean $\mathcal{F}(U) \to \mathcal{Q}(U)$ is a surjection for an open set $U$. 
\begin{nonexample}
    $X = \mathbb{P}^1, \mathcal{F} = \mathcal{O}_X$ and $I \subset \mathcal{O}_X$ the subsheaf consisting of sections vanishing at $0$ and $\infty$ then $\mathcal{O}_X \twoheadrightarrow \mathcal{O}_{X/I}$ but $$\C = \mathcal{O}_X(X) \to \mathcal{O}_{X/I}(X) = \C^2$$ is not surjective. 
\end{nonexample}
We now look towards closed embeddings and closed subschemes. Let $X = \Spec A$ be an affine scheme and $Y$ a scheme with $i : Y \to X$ a morphism. The following theorem illustrates the main idea given without proof. 
\begin{theorem}
    There exists an ideal $I \lhd A$ with $Y \cong \Spec(A/I)$ if and only if the following three conditions hold:
    \begin{enumerate}
        \item The image $i(Y) \subset X$ is closed
        \item $i : Y \to i(Y)$ is a homeomorphism
        \item $i^\sharp : \mathcal{O}_X \to i_*\mathcal{O}_Y$ is surjective
    \end{enumerate}
    In such a case we call $Y$ a closed subscheme of $X$ and $i$ a closed embedding.
\end{theorem}
So in general $i : Y \to X$ is a closed embedding if it is topologically (1 \& 2) and if $i^\sharp$ is surjective (3). The following are two non-examples of this idea. The first not fulfilling (2) and the second not fulfilling (3). 
\begin{nonexample}
    We consider the map $\A^1 \sqcup \A^1 \to \A^2$ where $t$ from the first $\A^1$ gets sent to $(t, 0)$ and $s$ from the second gets sent to $(s, 0)$. This map is not a homeomorphism onto its image as the origin of both $\A^1$ are sent to the origin of $\A^2$ and hence not injective. 
\end{nonexample}
\begin{nonexample}
    (A better non-example) We consider the map $i : \A^1_\C \to \A^2_\C$ where $a \mapsto (a^2, a^3)$. On the level of rings this is equivalent to the map $\C[x, y] \to \C[t]$ where $x \mapsto t^2$ and $y \mapsto t^3$. In this case the image of $i$ is closed and the the map from $\A^1_\C$ to its image is a homeomorphism. Now lets look at the map $i^\sharp : \mathcal{O}_{\A^2_\C} \to i_*\mathcal{O}_{\A^1_\C}$. On the level of stalks, this gives us the map $\C[x, y]_{(x, y)} \to \C[t]_{(t)}$ where $x \mapsto t^2$ and $y \mapsto t^3$. On this stalk however the map is not surjective as $t$ is not in the image. 
\end{nonexample}
We finish of this segment of the notes with the following equivalence stated without proof. 
\begin{proposition}
    TFAE
    \begin{enumerate}
        \item $i : Y \to X$ is a closed embedding
        \item For some open cover $\{U_i\}$ of $X$ we have that $i^{-1}(U_i) \to U_i$ are all closed embeddings
        \item For every open cover $\{U_i\}$ of $X$ we have that $i^{-1}(U_i) \to U_i$ are all closed embeddings
    \end{enumerate}
\end{proposition}

\newpage

\stepcounter{section}
\section*{Lecture 10: Morphisms to and from Projective Space}
Let \(k\) be a field.
Recall that we defined projective space \(\mathbf{P}^n_{k}\) such that the \(k\)-points of \(\mathbf{P}^n_{k}\) are 1-dimensional subspaces of \(k^{n+1}\).
A one-dimensional subspace is spanned by a non-zero vector, and two non-zero vectors span the same subspace if and only if they are related by scalar multiplication.
Given \(a_0,\dots,a_n \in k\), not all zero, we denote the subspace spanned by \((a_0,\dots,a_n)\) by \([a_0:\dots:a_n]\).
In this way, we can denote the \(k\)-points of \(\mathbf{P}^n_{k}\) by \([a_0:\dots:a_n]\).

Let us begin by constructing a map from \(\mathbf{P}^1_k\) to \(\mathbf{P}^2_k\) which, on $k$-points, sends \([a:b]\) to \([a^2:ab:b^2]\). We state without proof the following proposition, which will be very useful in this construction:
\begin{proposition}\label{prop:morphisms-form-a-sheaf}
If \(\{U_i\}\) is an open cover of a locally ringed space \(X\), and \(f_i:U_i\to Y\) are morphisms of locally ringed spaces such that \(f_i = f_j\) as morphisms from \(U_i\cap U_j\) to \(Y\), then there exists a morphism of locally ringed spaces \(f:X\to Y\). 
\end{proposition}
Now recall that \(\mathbf{P}^1_k=  \Spec(k[s])\cup \Spec(k[t])\), so we need maps from \(\Spec(k[s])\) and \(\Spec(k[t])\) to \(\mathbf{P}^2_k\). 

First, it suits us to rephrase our cover of \(\mathbf{P}^1_k\). Note that for any \(n\), we may take the cover of \(\mathbf{P}^n_k\) given by \(U_n = \Spec(k[\frac{x_0}{x_n},\frac{x_1}{x_n},\cdots,\frac{x_{n-1}}{x_n}])\) and so on, because the intersection \(U_{n-1}\cap U_n\) becomes 
\[U_{n-1}\cap U_n = \Spec\left(k\left[\frac{x_0}{x_n},\cdots,\frac{x_{n-1}}{x_n}\right]_{\frac{x_{n-1}}{x_n}}\right)\]
\[U_{n-1}\cap U_n = \Spec\left(k\left[\frac{x_0}{x_{n-1}},\cdots,\frac{x_{n}}{x_{n-1}}\right]_{\frac{x_n}{x_{n-1}}}\right)\]
These two rings are now exactly the same subring of \(k[x_0,\cdots,x_n]_{x_{n-1},x_n}\), so the gluing maps are now the identity. 

Returning to our map from \(\mathbf{P}^1_k\) to \(\mathbf{P}^2_k\) sending the \(k\)-points \([a:b]\) to \([a^2:ab:b^2]\), we use now the covers \(\mathbf{P}^1 = \Spec(k[\frac{a}{b}])\cup \Spec(k[\frac b a])\) and \(\mathbf{P}^2 = \Spec(k[\frac u w, \frac v w ])\cup \Spec(k[\frac u v,\frac w v])\cup \Spec(k[\frac v u, \frac w u])\). 

We finally define the map \(\Spec(k[\frac a b])\to \Spec(k[\frac u w,\frac v w])\) corresponding to the ring homomorphism with \(\frac u w \mapsto (\frac a b)^2\) and \(\frac v w\mapsto \frac a b\), and the map \(\Spec(k[\frac b a])\to \Spec(k[\frac v u,\frac w u])\) corresponding to the ring homomorphism with \(\frac v u \mapsto \frac b a\) and \(\frac w u \mapsto (\frac b a )^2\). On the overlaps, we have maps \(\Spec(k[\frac a b]_{\frac a b})\to \Spec(k[\frac u w,\frac v w]_{\frac u w})\) and \(\Spec(k[\frac b a ]_{\frac b a})\to \Spec(k[\frac v u, \frac w u]_{\frac w u})\) which are compatible with the gluing maps by our choices above. 

This defines our map \(i:\mathbf{P}^1_k\to \mathbf{P}^2_k\). We would like to check that \(i\) is a closed embedding; to do so, it suffices to check on a cover. For example, on the open set \(\Spec(k[\frac b a]) = U_0\), which is \(i^{-1}(V_0)\) for \(V_0 = \Spec(k[\frac w u, \frac v u])\), we see that \(i:\Spec(k[\frac b a])\to \Spec(k[\frac w u,\frac v u])\) arises from a surjective map of rings \(k[\frac w u,\frac v u]\to k[\frac b a]\), so \(U_0\cong \Spec(V_0/I)\): we have a closed embedding locally here (and similarly everywhere else). 

One can also obtain similar maps such as \(\mathbf{P}^1_k \to \mathbf{P}^3_k\) with \([a:b]\mapsto [a^3:a^2b:ab^2:b^3]\) and \(\mathbf{P}^2_k\to \mathbf{P}^5_k\) with \([x:y:z]\mapsto[x^2:xy:y^2:xz:z^2:yz]\). These are called Veronese embeddings, after Giuseppe Veronese (1854-1917). 


Now, what if we wanted a map with \([a:b]\mapsto[a+b:a-b]\)? We quickly see that the covers do not match as pleasantly as they did before. It would be much more useful to have a general idea of maps from any scheme \(X\) to \(\mathbf{P}^n\) as we did for affine schemes, where for affine \(Y\), maps \(X\to Y\) correspond to ``solving equations in \(\mathcal{O}_X(X)\)''. 

Consider \(\mathbf{P}^n_\mathbf{Z}\). A map \(\Spec (k)\to \mathbf{P}^n\) is a 1-dimensional subspace of \(k^{n+1}\). Could we likewise consider maps \(\Spec(R)\to \mathbf{P}^n\) as ``1-dimensional'' submodules of \(R^{n+1}\)? We will need the following:
\begin{definition}
For any locally ringed space \((X,\mathcal{O}_X)\), an \(\mathcal{O}_X\)-module \(\mathcal{M}\) is a sheaf of abelian groups with a multiplication \(\mathcal{O}_X(U)\times \mathcal{M}(U)\to \mathcal{M}(U)\) compatible with the sheaf restrictions.
\end{definition}
Just as a module in the traditional sense is free, we define what it means for an \(\mathcal{O}_X\)-module to be locally free:
\begin{definition}
If \((X,\mathcal{O}_X)\) a locally ringed space, \(\mathcal{M}\) an \(\mathcal{O}_X\)-module, we say \(\mathcal{M}\) is locally free of rank \(r\) if there exists a cover \(\{U_i\}\) of \(X\) such that \(\mathcal{M}|_{U_i} \cong \mathcal{O}_{U_i}^{\oplus r}\). 
\end{definition}
This allows us to classify morphisms from \(X\to \mathbf{P}^n\) with the following theorem (to be proven later?).
\begin{theorem}
    A morphism from \(X\) to \(\mathbf{P}^n\) is precisely a locally free \(\mathcal{O}_X\)-module of rank 1, \(\mathcal{L}\), and a locally split injection of \(\mathcal{O}_X\)-modules \(\mathcal{L}\to\mathcal{O}^{n+1}_X\). 
\end{theorem}
\begin{definition}
A split injection of \(O_X\)-modules is a map of \(O_X\)-modules \(f: \mathcal{L} \to \mathcal{M}\) such that there exists a map of \(O_X\)-modules \(g: \mathcal{M} \to \mathcal{L}\) such that \(g\circ f\) is the identity on \(\mathcal{L}\).
\end{definition}
We can define a split-injection more generally than for \(O_X\)-modules.
For example, if \(R\) is a ring, then we have an analogous notion of split injections of \(R\)-modules.
\begin{example}
  Let \(R = \mathbf{Z}\).
  The map of \(\mathbf{Z}\)-modules \(\mathbf{Z} \to \mathbf{Z}\) given by \(1 \mapsto 2\) is an injection, but not a split injection.
  The map \(\mathbf{Z}^2 \to \mathbf{Z}^3\) given by the matrix
  \[ \begin{pmatrix}
    1 & 2 \\
    2 & 1 \\
    3 & 0
   \end{pmatrix}
 \] is an injection but not a split injection.
 For a map \(\mathbf{Z}^r \to \mathbf{Z}^s\) given by an \(r \times s\) matrix \(M\) is split injection if and only if the determinants of the \(r \times r\) minors of \(M\) generate the unit ideal (Anand: I think this is true; and I am going to leave it as an exercise)!
\end{example}


%%% Local Variables:
%%% mode: LaTeX
%%% TeX-master: "main"
%%% End:

\stepcounter{section}
\newpage	
\section*{Lecture 11: Morphisms to Projective Space and Picard Group}

% Date: 1st September
% Note: ``Lecture 12'' (3rd September) was replaced with midsemester exam.

Let \(k\) be a field.
Previously, we have described morphisms $\Spec k \to \P^n_k$. We also briefly saw what morphisms $X \to {\P_\Z^n}$ might look like. In this lecture we will explore this idea further, and see how it can extend to morphisms $X \to \P^n$ for any scheme $X$. For this lecture, we use the shorthand $\P^n = \P_\Z^n$.

First, recall the case when $X = \Spec k$. Then a morphism $X \to \P^n$ is the same as the following data:
\begin{enumerate}
	\item A 1-dimensional $k$-vector space $\mathcal{L}$ (we use the letter $\mathcal{L}$ to remind ourselves that we are looking at one-dimensional vector spaces, or ``lines'').
	\item An injection $\mathcal{L} \hookrightarrow k^{n+1}$.
\end{enumerate}

To generalise different $X$, we need to consider how these requirements can be formulated for objects which are not as nice as vector spaces over fields. The construction that works is given by the following theorem.
\begin{theorem}
	A morphism of schemes $X \to \P^n$ is the same as the following data:
	\begin{enumerate}
		\item an $\mathcal{O}_X$-module $\mathcal{L}$ which is locally free of rank 1 (also called a ``line bundle''), and
		\item a locally split injection $\mathcal{L} \hookrightarrow \mathcal{O}_X^{n+1} = \bigoplus^{n+1} \mathcal{O}_X$
	\end{enumerate}
\end{theorem}

First, we should ask what a locally free, rank 1 $\mathcal{O}_X$-module actually looks like for a given scheme $X$. In fact, the set of such modules has a nice group structure called the Picard group.

\begin{definition}[Picard group]
	Given a scheme $X$, its \emph{Picard group} is denoted $\operatorname{Pic}(X)$, and is defined as the isomorphism classes:
	\begin{align*}
		\operatorname{Pic}(X) := \{ \mathcal{L} \; \vert \;  \text{$\mathcal{L}$ is a locally free, rank 1 $\mathcal{O}_X$-module.}  \}
	\end{align*}
	with multiplication given by the tensor product $\otimes$ of modules; the identity being the globally free module $\mathcal{O}_X$; and the inverse being the dual module $\mathcal{L}^* = \hom_{\mathcal{O}_X} (\mathcal{L}, \mathcal{O}_X)$, which satisfies $\mathcal{L} \otimes \mathcal{L}^* \cong \mathcal{O}_X$ in analogy to the vector space identity $V \otimes V^* \cong k$.
\end{definition}



\begin{example}
	We now see some examples of the Picard group.
	\begin{enumerate}
		\item Let $X = \Spec k = \{ \bullet \}$. A rank 1 $k$-module is just a 1 dimensional vector space, which is always globally free and unique up to isomorphism. So $\operatorname{Pic}(X) = \{\mathcal{O}_X\} \cong \{e\}$.
		\item In fact, whenever the underlying space $X = \{\bullet\}$ is a single point, then `local' and `global' are the same. So the only line bundle (locally free rank 1 module) is $\mathcal{O}_X$ itself, meaning $\operatorname{Pic}(X) = \{\mathcal{O}_X\} \cong \{e\}$. Geometrically, this corresponds to there being no room for a line to twist or change behaviour between open covers, because there is only one open cover.
		\item Let $R$ be a local ring, with the unique maximal ideal $m$. Then $X = \Spec R$ has a single closed point, $m$. Every non-empty open subset of $X$ must contain $m$. Because $m$ is in all parts of any open cover, the local properties around $m$ extend to global properties. So again, there is no distinction between `local' and `global' and $\operatorname{Pic}(X)$ is trivial.
		\item $\operatorname{Pic}(\Spec \Z) = \{\mathcal{O}_X\} \cong \{e\}$ (this is a consequence of the following theorem).
	\end{enumerate}
\end{example}

In fact, most of the affine schemes that we care about have trivial Picard group, due to the following theorem, which we state without proof.

\begin{theorem} \label{prop: pic trivial on ufd}
	Let $R$ be a ring which is a unique factorisation domain (for example, a PID), and let $X = \Spec R$. Then $\operatorname{Pic}(X) = \{\mathcal{O}_X\}$ is trivial.
\end{theorem}

\begin{remark}
	In number theory, the group $\operatorname{Pic}(\Spec R)$ is often called the \emph{class group} of $R$, and it measures the failure of unique factorisation in the ring $R$.
\end{remark}

We now know what (sheaves of) locally free, rank 1 $\mathcal{O}_X$-modules look like abstractly. But how do we construct them explicitly? 

\begin{definition}[Localisation of a module.]
	Let $R$ be a ring and $S \subseteq R$ be multiplicatively closed. The localisation of an $R$-module $M$ at $S$ is constructed the same way as the localisation of $R$ itself, namely:
	\begin{align*}
		S^{-1} M = \left\{ \frac{m}{s} \; \Big\vert \; m\in M, s \in S  \right\}
	\end{align*}
	with the appropriate equivalences.
\end{definition}

\begin{definition}[$\mathcal{O}_X$-module associated to a module.]
  Let $R$ be a ring and $M$ an $R$-module. Consider the affine scheme $X = \Spec R$, and write $D_f = \Spec R_f$. To each $D_f$, assign the localised module $M_f$, which is an $R_f$-module.  We can check that \(D_f \mapsto M_f\) defines a sheaf on the set of distinguished open, and hence extends uniquely to a sheaf on \(X\).
  We denote this extension by by $\Tilde{M}$.  It is a sheaf of $\mathcal{O}_X$-modules.

  For non-affine \(X\), we can construct \(O_X\)-modules by gluing together various \(\widetilde M\) on affine open patches and gluing.
\end{definition}

\begin{example}
	Some simple examples of this construction include:
	\begin{enumerate}
		\item If $M = R$, then $\Tilde{M} = \mathcal{O}_X$. 
		\item If $M = R \oplus \cdots \oplus R$, then $\Tilde{M} = \mathcal{O}_X \oplus \cdots \oplus \mathcal{O}_X$.
	\end{enumerate}
\end{example}

The tilde construction can also be used to see explicit examples of schemes with non-trivial Picard group.

\begin{example}
  Let $R = \Z[\sqrt{-5}]$ and let $X = \Spec R$.  Then \(R\) is not a UFD, and the failure of unique factorisation is where we will find a nontrivial element of $\operatorname{Pic}(X)$, due to \Cref{prop: pic trivial on ufd}.
  The non-trivial element of the Picard group will come from a place where unique factorisation fails:
  	\begin{align*}
		6 = 2 \cdot 3 = (1 + \sqrt{-5}) (1 - \sqrt{-5}).
	\end{align*}
	Consider the ideal $I = (2, 1 + \sqrt[]{-5}) \subset R$.
        We show that \(\widetilde I\) is locally free of rank 1, but \(I\) is not globally free of rank 1.
        To show local freeness, we take the open cover \(\Spec R = D_2 \cup D_3\).
        On $D_2$, the localisation $I_2 \subset R_2$ contains the unit 2, meaning we must have $1 \in I_2$ as well, and hence $I_2 = R_2$. This is a free module of rank 1. 
       On $D_3$, the localisation is $I_3 \subset R_3$. We can compute
	\begin{align*}
		\frac{1 - \sqrt{-5}}{3} (1 + \sqrt[]{-5}) = \frac{1 - (-5)}{3} = \frac{6}{3} = 2
	\end{align*}
	Since $\frac{1 - \sqrt{-5}}{3} \in R_3$, this shows that 2 is a multiple of $1 + \sqrt{-5}$ in $R_3$, and hence $I_3 = (1 + \sqrt{-5})$. Since $1 + \sqrt{-5}$ is not a zero divisor, $I_3$ is a free module of rank 1.
        Therefore, \(I\) is locally free of rank 1.

        To show that \(I\) not globally free, we need to show that \(I \subset R\) is not a principal ideal.
        Since \(2\) and \(1+\sqrt{-5}\) are both irreducible in \(R\), they have no non-unit common factor.
        But \(I\) is not the unit ideal; so it can't be principal. 
\end{example}


\begin{nonexample}
	For $X = \Spec \C[x,y]$, consider the ideal $I = (x,y)$. Its corresponding $\mathcal{O}_X$-module $\Tilde{I}$ is not locally free of rank 1.
\end{nonexample}


\begin{proposition}
	If an $\mathcal{O}_X$ module $\mathcal{F}$ is locally free of rank $r$, then all stalks are free \(\mathcal{O}_{X,x}\)-modules of rank $r$.  That is, for all $x \in X$, we have isomorphisms $\mathcal{O}_{X,x}^{\oplus r} \cong \mathcal{F}_x$.
\end{proposition}
\begin{proof}
	If $\mathcal{F}$ is a locally free $\mathcal{O}_X$-module of rank $r$, then there is an open cover of $X$ such that for each open $U$, we have $\mathcal{F}|_U \cong \mathcal{O}_U^{\oplus r}$. Taking stalks is compatible with direct sums, so for any point $x \in X$, this local isomorphism becomes $\mathcal{F}_x \cong \left( \mathcal{O}_X^{\oplus r} \right)_x \cong \mathcal{O}_{X,x}^{\oplus r}$, as an $\mathcal{O}_X$-module isomorphism. 
\end{proof}

\begin{proposition}
	\begin{align*}
		\operatorname{Pic}(\P^n_k) = \Z
	\end{align*}
\end{proposition}

\begin{proof}[Proof for $n=1$.]
	We can construct the locally free $\mathcal{O}_X$-modules of rank 1 for $X = \P^1_k$, which are the elements of the Picard group. 
	
	$X$ is open covered by $U_1 = \Spec k[t]$ and $U_2 = \Spec k[s]$ with transition map $t = s^{-1}$. On $U_1$, we have $\mathcal{L}$ being a free $k[t]$-module of rank 1, generated by some $e_1$. And on $U_2$, it is a free $k[t^{-1}]$-module of rank 1, generated by some $e_2$. 
	
	On the intersection, we need an isomorphism $\varphi: k[t,t^{-1}] \cdot e_1 \to k[t, t^{-1}] \cdot e_2$ of $k[t, t^{-1}]$-modules. If we make the forwards map be $\varphi(e_0) = f(t) e_1$ for some $f \in k[t, t^{-1}]$, then in order to have an inverse map, we require $f$ to be a unit, so that there is some $g$ such that $f(t) g(t) = 1$. With this, the inverse is $\varphi(e_1) = g(t) e_0$.
	
	Our choices for the isomorphism are entirely determined by the choice $f \in k[t, t^{-1}]$. The only units (and hence the only valid choices) are nonzero monomials $f(t) = a \cdot t^n$ for some $a \in k^\times$ and $n \in \Z$.
	
	Notice that the choice of $a \in k^\times$ does not matter, up to sheaf isomorphism. Consider constructing two sheaves: $\mathcal{F}$ using the isomorphism given by $at^n$, and $\mathcal{G}$ using the isomorphism given by $t^n$. Then the following diagram commutes, so $\mathcal{F} \cong \mathcal{G}$.
	% https://q.uiver.app/#q=WzAsNCxbMCwwLCJcXG1hdGhjYWx7Rn0oVV8xKSJdLFswLDIsIlxcbWF0aGNhbHtGfShVXzIpIl0sWzIsMiwiXFxtYXRoY2Fse0d9KFVfMikiXSxbMiwwLCJcXG1hdGhjYWx7R30oVV8xKSJdLFswLDMsIjEiXSxbMSwyLCJhXnstMX0iXSxbMCwxLCJhIHRebiIsMV0sWzMsMiwidF5uIiwxXV0=
	\[\begin{tikzcd}
		{\mathcal{F}(U_1)} && {\mathcal{G}(U_1)} \\
		\\
		{\mathcal{F}(U_2)} && {\mathcal{G}(U_2)}
		\arrow["{\times 1}", from=1-1, to=1-3]
		\arrow["{a t^n}"{description}, from=1-1, to=3-1]
		\arrow["{t^n}"{description}, from=1-3, to=3-3]
		\arrow["{\times a^{-1}}", from=3-1, to=3-3]
	\end{tikzcd}\]
	
	However, the choice of $n \in \Z$ \textit{does} matter, because it is not possible to glue different powers of $t$ in a way that retains the sheaf structure. 
	
	So up to sheaf isomorphism, there are $n \in \Z$ choices of the $k[t,t^{-1}]$-module isomorphism, and hence an element of the Picard group corresponding to every such $n$. Thus $\operatorname{Pic}(X) = \Z$. 
      \end{proof}

      \begin{remark}
        We will later introduce the Cech complex that will make the procedure above more systematic.
      \end{remark}



%%% Local Variables:
%%% mode: LaTeX
%%% TeX-master: "main"
%%% End:

\newpage

\stepcounter{section}
\section*{Lecture 13: Line bundles and the Picard group}
Recall that morphisms from a locally ringed space $(X, \mathcal{O}_X)$ to projective space $\P^n$ are in bijection with locally free $\mathcal{O}_X$-modules $\mathcal{L}$ of rank 1 (also known as line bundles), along with a locally split injection of $\mathcal{O}_X$-modules $\mathcal{L} \hookrightarrow \mathcal{O}_X^{n+1}$, up to isomorphism. In this context, an isomorphism between $\mathcal{L}_1 \hookrightarrow \mathcal{O}_X^{n+1}$ and $\mathcal{L}_2 \hookrightarrow \mathcal{O}_X^{n+1}$ is an $\mathcal{O}_X$-module isomorphism between $\mathcal{L}_1$ and $\mathcal{L}_2$ such that the following diagram commutes: 
\[
\begin{tikzcd}
    \mathcal{L}_1  && \mathcal{O}_X^{n+1} \\
    \\
    \mathcal{L}_2 && \mathcal{O}_X^{n+1}.
    \arrow[hook, from=1-1, to=1-3]
    \arrow[hook, from=3-1, to=3-3]
    \arrow["\cong"{description}, from=1-1, to=3-1]
    \arrow[shift right, no head, from=1-3, to=3-3]
	\arrow[shift left, no head, from=1-3, to=3-3]
\end{tikzcd}
\]

We now look at maps between line bundles. Let $\mathrm{Hom}_{\mathcal{O}_X}(\mathcal{L}_1, \mathcal{L}_2)$ denote the set of $\mathcal{O}_X$-linear maps of sheaves from $\mathcal{L}_1$ to $\mathcal{L}_2$. The simplest case is when both $\mathcal{L}_1$ and $\mathcal{L}_2$ are just $\mathcal{O}_X$.
\begin{proposition}
  We have a natural bijection
  \[ \mathrm{Hom}_{\mathcal{O}_X}(\mathcal{O}_X, \mathcal{O}_X) \simeq \mathcal{O}_X(X)\]
  defined as follows.
  To go from the left to the right, we send a homomorphism \(\phi\) to \(\phi(1)\).
  To go from the right to the left, we send \(f \in \mathcal{O}_X(X)\) to the homomorphism given by multiplication by \(f\).
\end{proposition}
\begin{proof}
  It is clear that \(\mathcal{O}_X(X) \to \mathrm{Hom}_{\mathcal{O}_X}(\mathcal{O}_X, \mathcal{O}_X) \to \mathcal{O}_X(X)\) is the identity.
  Let us show that \(\mathrm{Hom}_{\mathcal{O}_X}(\mathcal{O}_X, \mathcal{O}_X) \to \mathcal{O}_X(X) \to \mathrm{Hom}_{\mathcal{O}_X}(\mathcal{O}_X, \mathcal{O}_X)\) is the identity.
  Let \(\phi \in \mathrm{Hom}_{\mathcal{O}_X}(\mathcal{O}_X, \mathcal{O}_X)\) be given and let \(f = \phi(1)\).
  On any \(U \subset X\), consider 
  \[\phi(U) \colon \mathcal{O}_X(U) \to \mathcal{O}_X(U).\]
  Since \(\phi\) is a map of sheaves, it must sends the restriction of \(1\) to the restriction of \(f\).
  By \(\mathcal{O}_X\)-linearity, it must be multiplication by (the restriction of) \(f\).
\end{proof}

By similar reasoning, have
\[ \mathrm{Hom}_{\mathcal{O}_X}(\mathcal{O}_X, \mathcal{L}) \simeq \mathcal{L}(X).\]
We now consider $\mathrm{Hom}_{\mathcal{O}_X}(\mathcal{L}, \mathcal{L})$.
\begin{proposition}\label{prop:homll}
  The map
  \[ \mathcal{O}_X(X) \to \mathrm{Hom}_{\mathcal{O}_X}(\mathcal{L}, \mathcal{L})\]
  that sends \(f\) to the homomorphism given by multiplication by \(f\) is a bijection.
\end{proposition}
\begin{proof}
  Let $\{U_i\}$ be an open cover of $X$ such that $\mathcal{L}|_{U_i} \cong O_{U_i}$ for each $i$, also known as a trivialisation of $\mathcal{L}$. A map $\alpha: \mathcal{L} \to \mathcal{L}$ restricts to maps $\alpha_i: \mathcal{L}|_{U_i} \to \mathcal{L}|_{U_i}$. Furthermore, on overlaps $U_{ij} := U_i \cap U_j$, we see that the restrictions of $\alpha_i$ and $\alpha_j$ to $U_{ij}$ must agree. As $\mathcal{L}|_{U_i} \cong \mathcal{O}_{U_i}$, we see that each $\alpha_i$ is multiplication by an element of $\mathcal{O}_X(U_i)$, and by our ``gluing condition" on $\alpha_i$'s we know that the multiplication has to be by an element that agrees on overlaps. Hence every element of $\mathrm{Hom}_{\mathcal{O}_X}(\mathcal{L}, \mathcal{L})$ is just multiplication by a $\lambda \in \mathcal{O}_X(X)$, and we have $\mathrm{Hom}_{\mathcal{O}_X}(\mathcal{L}, \mathcal{L}) \simeq \mathcal{O}_X(X)$
\end{proof}
A consequence of \Cref{prop:homll} is that the set of isomorphisms from \(L\) to \(L\) is the set of invertible elements of \(\mathcal{O}_X(X)\):
\[ \mathrm{Iso}_{\mathcal{O}_X}(\mathcal{L}, \mathcal{L}) \simeq \mathcal{O}_X(X)^\times.\]


\begin{remark}
    Given $\{\alpha_i\}$ such that the restrictions agree on overlaps, there exists a unique $\alpha$ such that $\alpha$ restricted to each $U_i$ is precisely $\alpha_i$. In other words, given sheaves $\mathcal{F}, \mathcal{G}$ on $X$, $U \mapsto \mathrm{Hom}(\mathcal{F}|_U, \mathcal{G}|_U)$ is also a sheaf. 
\end{remark}

Recall that the Picard group of $X$, denoted $\mathrm{Pic}(X)$, is the group of isomorphism classes of line bundles on $X$. Fix an open covering $\{U_i\}$ of $X$, and let $\mathrm{Pic}^{\{U_i\}}(X)$ denote the subgroup of line bundles trivialised by $\{U_i\}$. For some $\mathcal{L} \in \mathrm{Pic}^{\{U_i\}}(X)$, we choose isomorphisms $\alpha_i: \mathcal{O}_{U_i} \to \mathcal{L}|_{U_i}$. Then by above, the restrictions of $\alpha_i$ and $\alpha_j$ to $U_{ij}$ differ by an element of $\mathcal{O}_X(U_{ij})^\times$; that is, $\alpha_j = \lambda_{ij} \cdot \alpha_i$, where $\lambda_{ij} \in \mathcal{O}_X(U_{ij})^\times$. On triple overlaps $U_{ijk}$, the previous identity then tells us that $\lambda_{ij} \cdot \lambda_{jk} \cdot \lambda_{ki} = 1$. 

So far we have characterised $\mathcal{L}$ by a choice of isomorphisms $\alpha_i: \mathcal{O}_{U_i} \to \mathcal{L}|_{U_i}$ along with elements $\lambda_{ij} \in \mathcal{O}_X(U_{ij})$ satisfying $\lambda_{ij} \cdot \lambda_{jk} \cdot \lambda_{ki} = 1$. What happens if we change $\alpha_i$? Let $\alpha_i'$ be another isomorphism between $\mathcal{O}_{U_i}$ and $\mathcal{L}|_{U_i}$. We know that $\alpha_i' = \alpha_i \cdot \mu_i$ for some $\mu_i \in \mathcal{O}_X(U_i)^\times$, which in turn tells us that $\lambda_{ij}' = \lambda \cdot \mu_j\mu_i^{-1}$. This motivates us to construct the following maps between abelian groups:
\[
\begin{tikzcd}
    \overbrace{\prod_i \mathcal{O}_X(U_i)^\times}^{C_0} && \overbrace{\prod_{i<j} \mathcal{O}_X(U_{ij})^\times}^{C_1} && \overbrace{\prod_{i<j<k} \mathcal{O}_X(U_{ijk})^\times}^{C_2} \\
    (\mu_i) && (\mu_j \mu_i^{-1}) && \\
    && (\lambda_{ij}) && (\lambda_{ij} \cdot \lambda_{jk} \cdot \lambda_{ik}^{-1})
    \arrow[maps to, from=2-1, to=2-3]
	\arrow[maps to, from=3-3, to=3-5]
    \arrow[from=1-1, to=1-3]
	\arrow[from=1-3, to=1-5]
\end{tikzcd}
\]
We note that the image of the first map lies in the kernel of the second map, so we have a (co)chain complex. Now, a line bundle $\mathcal{L} \in \mathrm{Pic}^{\{U_i\}}(X)$ can be characterised by an element of $\ker(C_1 \to C_2)$ (representing the elements $\lambda_{ij} \in \mathcal{O}_X(U_{ij})$ satisfying $\lambda_{ij} \cdot \lambda_{jk} \cdot \lambda_{ki} = 1$) modulo the image of $C_0 \to C_1$ (as the choices of isomorphisms $\alpha_i: \mathcal{O}_{U_i} \to \mathcal{L}|_{U_i}$ should not matter). Hence $\mathrm{Pic}^{\{U_i\}}(X)$ is isormophic to $\ker(C_1 \to C_2)/\mathrm{im}(C_0 \to C_1) = H^1(C_*(\{U_i\}))$, the first cohomology group of our complex. 

For a refinement $\{V_i\}$ of $\{U_i\}$, we have the inclusion $\mathrm{Pic}^{\{U_i\}}(X) \subset \mathrm{Pic}^{\{V_i\}}(X)$, which induces a map $H^1(C_*(\{U_i\})) \to  H^1(C_*(\{V_i\}))$. Thus we may take 
\[
\mathrm{Pic}(X) = \varinjlim_{\mathrm{refinement}} H^1(C_*(\{U_i\}))
\]
as a definition of the Picard group. Furthermore, if we can find an open cover $\{U_i\}$ of $X$ such that $\mathrm{Pic}(U_i)$ is trivial for each $i$, then we can claim that $\mathrm{Pic}(X) = \mathrm{Pic}^{\{U_i\}}(X)$, as further refinements will not change $\mathrm{Pic}^{\{U_i\}}(X)$. 
\begin{example}
    We compute the Picard group of projective space $\P^n_k$ for a field $k$, using the standard cover of $\P^n_k$ by affine spaces (since $\mathrm{Pic}(\A^n_k)$ is trivial).

    For $n=1$, as there are no triple intersections, $C_2 = \{1\}$. Then $C_1$ are the units of $k[t, t^{-1}]$, which is given by $\{a \cdot t^n \mid a \in k^\times, n \in \Z\}$. Finally, the units of $k[x]$ are just $k^\times$, so $C_0 = k^\times \times k^\times$. Hence we have
    \[
    \begin{tikzcd}
    k^\times \times k^\times && \{a \cdot t^n \mid a \in k^\times, n \in \Z\} \cong k^\times \times (\Z, +) && \{1\} \\
    (\lambda_1, \lambda_2) && \lambda_2 \lambda_1^{-1} \cdot t^0 && 
    \arrow[maps to, from=2-1, to=2-3]
    \arrow[from=1-1, to=1-3]
	\arrow[from=1-3, to=1-5]
    \end{tikzcd}
    \]
    As everything in $C_1$ is in the kernel, we see that $\mathrm{ker}/\mathrm{im} \cong \Z$. 

    For $n = 2$, for clarity of notation let $U_1 = \Spec k[\frac{x_2}{x_1}, \frac{x_3}{x_1}]$, $U_2 = \Spec k[\frac{x_1}{x_2}, \frac{x_3}{x_2}]$ and $U_3 = \Spec k[\frac{x_1}{x_3}, \frac{x_2}{x_3}]$. Then $U_1 \cap U_2$ can be seen as $\Spec \, k[\frac{x_2}{x_1}, \frac{x_3}{x_1}, \frac{x_1}{x_2}]$, so as before the units are given by $\{a \cdot (\frac{x_1}{x_2})^n\}$. In the triple overlap, all variables become invertible, so the units are of the form $\{a \cdot (\frac{x_1}{x_2})^n(\frac{x_1}{x_3})^m\}$. Then we have
    \[
    \begin{tikzcd}
    (k^\times)^3 & \{a \cdot (\frac{x_1}{x_2})^n\} \times \{b \cdot (\frac{x_2}{x_3})^m\} \times \{c \cdot (\frac{x_3}{x_1})^l\} & \{a \cdot (\frac{x_1}{x_2})^n(\frac{x_1}{x_3})^m\}
    \arrow[from=1-1, to=1-2]
	\arrow[from=1-2, to=1-3]
    \end{tikzcd}
    \]
    For an element of $C_1$ to be in the kernel, we must have $(\frac{x_1}{x_2})^n(\frac{x_2}{x_3})^m(\frac{x_3}{x_1})^l = 1$, as well as $abc = 1$. A choice of $n$ determines both $m$ and $l$, so the Picard group is once again $\Z$, since changing $a, b, c$ is equivalent to changing by some element in the image from $C_0$.

    Larger values of $n$ follow similarly to $\P^2_k$.  
\end{example}
%%% Local Variables:
%%% mode: LaTeX
%%% TeX-master: "main"
%%% End:

\newpage

\stepcounter{section}
\section*{Lecture 14: Products}
Let $X,Y$ be schemes. In the same way that we can take the product of two sets or groups we want to define a product scheme $X\times Y$. This could mean taking the set of points $X\times Y$ and putting a scheme structure on it, i.e. topologising the set and providing a sheaf of rings. However this turns out to not be ``the right thing'' to do.
To see why, we must first pin down what we mean by ``the right thing'' in this context.

To determine the correct definition we must specify the properties we want the product to satisfy. Our intuition tells us that we should have maps $\pi_X\colon X\times Y \to X$ and $\pi_Y\colon X\times Y \to Y$ such that for any scheme $Z$ we have
\[
\text{Maps}(Z, X\times Y)\to \text{Maps}(Z,X) \times \text{Maps}(Z,Y),
\]
\[
\varphi\mapsto (\pi_X\circ \varphi, \pi_Y\circ \varphi),
\]
and this assignment is a bijection. This leads us to define the product as the triple $(X\times Y, \pi_X, \pi_Y)$ that satisfies the following commutative diagram.
\[\begin{tikzcd}
	&&& X \\
	Z && {X\times Y} \\
	&&& Y
	\arrow[curve={height=-6pt}, from=2-1, to=1-4]
	\arrow["{\exists!}", dashed, from=2-1, to=2-3]
	\arrow[curve={height=6pt}, from=2-1, to=3-4]
	\arrow["{\pi_X}"', from=2-3, to=1-4]
	\arrow["{\pi_Y}", from=2-3, to=3-4]
\end{tikzcd}\]
We then ask two questions,
\begin{enumerate}
    \item Does such an object exist?
    \item Is such an object unique?
\end{enumerate}
The answer to the second one is typically false, as any object satisfying the diagram can be relabelled, so we instead ask if such an object is unique up to unique isomorphism. That is, if $\square$ and $\bigcirc$ are two such specified objects there exists a unique isomorphism $\square\to\bigcirc$ compatible with the projection maps.

Suppose $\square$ and $\bigcirc$ are two objects that fit the commutative diagram. Then we can let $Z = \square$ to obtain that the identity 
id$_\square$ is the unique map which makes the diagram commute. We then use the fact that $\bigcirc$ is also a product to obtain the following diagram,
\[\begin{tikzcd}
	&&&&& X \\
	\square && \bigcirc && \square \\
	&&&&& Y
	\arrow[curve={height=-18pt}, from=2-1, to=1-6]
	\arrow["{\exists! f}", dashed, from=2-1, to=2-3]
	\arrow[curve={height=18pt}, from=2-1, to=3-6]
	\arrow[curve={height=-6pt}, from=2-3, to=1-6]
	\arrow["{\exists! g}", dashed, from=2-3, to=2-5]
	\arrow[curve={height=6pt}, from=2-3, to=3-6]
	\arrow[from=2-5, to=1-6]
	\arrow[from=2-5, to=3-6]
\end{tikzcd}\]
so $g\circ f$ is also a map which works. By uniqueness of the map we obtain Id$_\square = g\circ f$, and similarly by swapping $\square$ and $\bigcirc$ we obtain Id$_\bigcirc = f\circ g$. Thus $f$ is the unique isomorphism from $\square$ to $\bigcirc$ as desired.

We have shown that if a product exists then it is unique up to unique isomorphism, so we denote the product by $X\times Y$. We now consider the question of existence.

\begin{example}
(Sets)
Let $X,Y$ be sets, we want to define the product $(X\times Y, \pi_X, \pi_Y)$. We claim that the Cartesian product
\[
X\times Y \colonequals \{ (x,y) \mid x\in X, y\in Y\},
\]
with $\pi_X,\pi_Y$ the projection maps to the first and second coordinate respectively, satisfies the specification. For any test set $T$ with maps $f,g$ we obtain the diagram,
\[\begin{tikzcd}
	&&& X \\
	T && {X\times Y} \\
	&&& Y
	\arrow["f", curve={height=-30pt}, from=2-1, to=1-4]
	\arrow["{t\mapsto(f(t),g(t))}", from=2-1, to=2-3]
	\arrow["g"', curve={height=30pt}, from=2-1, to=3-4]
	\arrow[from=2-3, to=1-4]
	\arrow[from=2-3, to=3-4]
\end{tikzcd}\]
where the uniqueness of the middle arrow is because for any $\varphi\colon T \to X\times Y$ satisfying the diagram we must have $\pi_X\circ \varphi = f$ and $\pi_Y\circ \varphi = g$, so $\varphi(t) = (f(t),g(t))$.
\end{example}

\begin{example}\hfill
\begin{itemize}
    \item (Topological spaces) Let $X,Y$ be topological spaces. Then $X\times Y$ is the Cartesian product with the product topology.
    \item (Groups) Let $G,H$ be groups. Then $G\times H$ is the Cartesian product with component-wise multiplication.
\end{itemize}
\end{example}

In all the examples we have done thus far the product is the Cartesian product with some additional structure imposed on it.
Surprisingly, this is not the case with schemes.
Here is the first example.

\begin{example}
    Let $X=Y=\Spec \Z$. Then the data of a test scheme $Z$ with maps $f\colon Z\to X$, $g\colon Z\to Y$ is equivalent to ring maps $\varphi\colon \Z\to \mathcal{O}_Z(Z)$, $\rho\colon \Z\to \mathcal{O}_Z(Z)$. However there is a unique ring map from $\Z$ to any ring, so $\rho=\varphi$ and the data of the map from $Z$ to $X\times Y$ is equivalent to the ring map $\varphi$. Thus we obtain that $X\times Y = \Spec \Z$, which satisfies the commutative diagram
    \[\begin{tikzcd}
	&&& {\Spec \Z} \\
	Z && {\Spec \Z} \\
	&&& {\Spec \Z}
	\arrow["f", from=2-1, to=1-4]
	\arrow[dashed, from=2-1, to=2-3]
	\arrow["g"', from=2-1, to=3-4]
	\arrow[from=2-3, to=1-4]
	\arrow[from=2-3, to=3-4]
\end{tikzcd}\]
\end{example}

\begin{example}
    Let $X = \Spec (\Z/2\Z)$ and $Y = \Spec (\Z/3\Z)$. Then $X\times Y = \emptyset$ is the empty scheme. This is because a scheme $Z$ with maps $f\colon Z\to X$, $g\colon Z\to Y$ is equivalent to ring maps $\varphi\colon \F_2\to \mathcal{O}_Z(Z)$, $\rho\colon \F_3\to \mathcal{O}_Z(Z)$. However we observe that
    \[
    1_Z = 3\rho(1) - 2\varphi(1) = \rho(3) - \varphi(2) = 0_Z,
    \]
    so $\mathcal{O}_Z(Z)$ must be the zero ring, and therefore \(Z\) must be the empty scheme.
    In other words, the only scheme that maps to \(X\) and \(Y\) is the empty scheme.
    As a result, the product \(X \times Y\) must also be the empty scheme, so we obtain the diagram
    \[\begin{tikzcd}
	&&& {\Spec \F_2} \\
	\emptyset && \emptyset \\
	&&& {\Spec \F_3}
	\arrow["f", from=2-1, to=1-4]
	\arrow[dashed, from=2-1, to=2-3]
	\arrow["g"', from=2-1, to=3-4]
	\arrow[from=2-3, to=1-4]
	\arrow[from=2-3, to=3-4]
\end{tikzcd}\]
\end{example}

\begin{example}
    Let $X = \Spec \Z[x] $ and $Y = \Spec \Z[y]$ be two copies of $\A^1_\Z $. Then $X\times Y = \Spec \Z[x,y] = \A_\Z^2$. To see why let $T$ be a test scheme with maps $ T\to X$, $T\to Y$, then this data is equivalent to ring maps $\varphi\colon \Z[x] \to \mathcal{O}_T(T)$, $\rho\colon \Z[y]\to \mathcal{O}_T(T)$. However the data of $(\varphi,\rho)$ is the same as choosing $x\mapsto f\in \mathcal{O}_T(T)$ and $y\mapsto g\in \mathcal{O}_T(T)$, which uniquely specifies a map $\Z[x,y] \to\mathcal{O}_T(T)$. All this information is then captured by the following diagram.
    \[\begin{tikzcd}
	&&& {\Spec \Z[x]} \\
	T && {\Spec \Z[x,y]} \\
	&&& {\Spec \Z[y]}
	\arrow["{x\mapsto f}", curve={height=-24pt}, from=2-1, to=1-4]
	\arrow["{(x,y)\mapsto(f,g)}", dashed, from=2-1, to=2-3]
	\arrow["{y\mapsto g}"', curve={height=24pt}, from=2-1, to=3-4]
	\arrow[from=2-3, to=1-4]
	\arrow[from=2-3, to=3-4]
\end{tikzcd}\]

By the same argument we also have for $\A^n_\Z = \Spec \Z[x_1,\dots, x_n]$ and $\A^m_\Z = \Spec \Z[y_1,\dots, x_m]$ that $\A^n_\Z\times \A^m_\Z = \A^{n+m}_\Z$.
\end{example}

We now note that the set of prime ideals of \(\Z[x,y]\) is not the product of the set of prime ideals of \(\Z[x]\) and \(\Z[y]\).
That is, we do not have the equality of sets
\[
\Spec \Z[x,y] = \Spec \Z[x] \times \Spec \Z[y].
\]

All the examples above contradict our intuition that the points of a product should be the product of the points of the parts.
Nevertheless, the intuition is correct if we use a different notion of points than the literal points of the underlying topological space.

We recall the two notions of points of a scheme $X$:
\begin{enumerate}
    \item The actual points of $X$ as a topological space,
    \item Given a field (or even a ring) \(k\), the $k$-points of \(X\), namely
      \[ \operatorname{Maps}(\Spec k, X).\]
\end{enumerate}
We often think of the first type of points as an implementation detail, defined purely to get a topology that admits a good sheaf of rings.
The second notion of points is more fundamental.
For example, suppose \(X\) is the affine scheme \(\Spec \Z[x_1,\dots,x_n]/(f_1,\dots,f_m)\).
Then the set of \(k\)-points of \(X\) is precisely the set of \(k\)-valued solutions to the system of equations
\[ f_1(x_1,\dots,x_n) = 0, \dots, f_m(x_1,\dots,x_n) = 0.\]
The primary concern of algebraic geometry is understanding solutions of polynomial equations!

The definition of the product scheme \(X \times Y\) behaves exactly as intended for this second, more fundamental notion of a point.
By definition, we have
\[ \operatorname{Maps}(\Spec k, X \times Y) = \operatorname{Maps}(\Spec k, X) \times \operatorname{Maps}(\Spec k,Y). \]
For any \(k\), the \(k\)-valued points of \(X \times Y\) are precisely the product of the \(k\)-valued points of \(X\) and \(Y\).

We now turn our attention to the general question: given \(X\) and \(Y\), does \(X \times Y\) exist?
Let us take the simpler case where both \(X\) and \(Y\) are affine, say $\Spec A$ and $Y =\Spec B$.
Then asking for \(X \times Y\) is asking whether there is a scheme $\square$ which satisfies the universal property in the following diagram: 
\[\begin{tikzcd}
	&&& {\Spec A} \\
	T && \square \\
	&&& {\Spec B}
	\arrow["{A\to \mathcal{O}_T(T)}", curve={height=-0pt}, from=2-1, to=1-4]
	\arrow[dashed, from=2-1, to=2-3]
	\arrow["{B\to \mathcal{O}_T(T)}"', curve={height=0pt}, from=2-1, to=3-4]
	\arrow[from=2-3, to=1-4]
	\arrow[from=2-3, to=3-4]
      \end{tikzcd}.\]
    We can try to find an affine scheme \(\square\) that satisfies the property.
    Trying to find an affine \(\square\) is equivalent to trying to find a ring $R$ with maps $A\to R$ and $B\to R$ satisfying the following universal property:
\[\begin{tikzcd}
	&&& A \\
	S && R \\
	&&& B
	\arrow[from=1-4, to=2-1]
	\arrow[from=1-4, to=2-3]
	\arrow[dashed, from=2-3, to=2-1]
	\arrow[from=3-4, to=2-1]
	\arrow[from=3-4, to=2-3]
\end{tikzcd}\]

Such a ring does exist and we will now construct it.
Consider the set 
\[
G \colonequals \{ a\otimes b \mid a\in A, b\in B\}.
\]
Let \(A \otimes B\) be the free abelian group generated by \(G\) modulo the following relations:
\[
  \begin{split}
    a_1 \otimes b + a_2\otimes b - (a_1 +a_2)\otimes b, \text{ for all } a_1, a_2 \in A \text{ and } b \in B,\\
    a\otimes b_1 + a\otimes b_2 - a \otimes (b_1 + b_2), \text{ for all } a \in A \text{ and } b_1,b_2 \in B.
  \end{split}
\]
We make \(A \otimes B\) a ring by setting
\[ (a_1 \otimes b_1) \times (a_2 \otimes b_2) = a_1a_2 \otimes b_1b_2,\]
and extending by linearity.
We need to check that this is well-defined modulo the relations, and satisfies the ring axioms, but we skip those checks.
We define maps $A \to A\otimes B$ and $B \to A\otimes B$ by $a\mapsto a\otimes 1$ and $b\mapsto 1\otimes b$ respectively.
Again, it is easy to check that these are ring homomorphisms.

Given maps $f\colon A \to S$ and $g\colon B \to S$, there is a unique map \(h \colon A \otimes B \to S\) such that the composite \(A \to A \otimes B \to S\) is \(f\) and \(B \to A \otimes B \to S\) is \(g\).
This map is given by
\[ h(a \otimes b) = f(a) g(b).\]
It is easy to check that this is a well-defined homomorphism.

We conclude that $\Spec A \times \Spec B$ exists and it equal to $\Spec(A\otimes B)$. As an exercise show that $\F_2\otimes \F_3 = 0$, which agrees with what we expected from before.
%%% Local Variables:
%%% mode: LaTeX
%%% TeX-master: "main"
%%% End:

\newpage

\stepcounter{section}
\section*{Lecture 15: Fibered Products}
In this lecture we turn our attention to a categorical construction which will be a sort of generalisation of the product. \\
The setup will consist of three schemes $X,Y,Z$, together with morphisms $f:X\to Z,g:Y\to Z$. The fibered product will be a universal object that is the `closest' to this configuration, in the following sense:
\begin{definition}[Fibered product]
    Let $X,Y,Z$ be schemes, and $f:X\to Z,g:Y\to Z$ be morphisms. A fibered product of this configuration is a scheme $\square$, together with morphisms $p_X:\square\to X,p_Y:\square\to Y$ which make the following diagram commute:
\[\begin{tikzcd}
	\square & X \\
	Y & Z
	\arrow["{p_X}", from=1-1, to=1-2]
	\arrow["{p_Y}"', from=1-1, to=2-1]
	\arrow["f", from=1-2, to=2-2]
	\arrow["g"', from=2-1, to=2-2]
\end{tikzcd}\]
    Moreover, a fibered product must satisfy the following universal property. Whenever $T$ is a scheme with maps $T\to X,T\to Y$ which commute with $f,g$, we have a unique morphism $\varphi:T\to\square$ which makes the following diagram commute:
\[\begin{tikzcd}
	T && \\
	& \square & X \\
	& Y & Z
	\arrow["{\exists!\varphi}"{description}, dashed, from=1-1, to=2-2]
	\arrow[curve={height=-12pt}, from=1-1, to=2-3]
	\arrow[curve={height=12pt}, from=1-1, to=3-2]
	\arrow["{p_X}", from=2-2, to=2-3]
	\arrow["{p_Y}"', from=2-2, to=3-2]
	\arrow["f", from=2-3, to=3-3]
	\arrow["g"', from=3-2, to=3-3]
\end{tikzcd}\]
\end{definition}
If a fibered product exists, it is unique up to unique isomorphism by the same argument as for the product, and so we refer to it as \textit{the} fibered product, and denote it by $X\times_ZY$. As is often the case in mathematics, this notation omits the most important information: the morphisms involved. \\
Of course this definition need not be restricted to schemes, and works for any category. Let's go through some examples in more familiar categories.
\begin{example}[Sets]
    Let $X,Y,Z$ be sets, and $f:X\to Y,g:X\to Z$ functions. The fibered product, if it exists, is a set with maps to $X$ and $Y$, and so by the universal property of the product, there is a unique map $\varphi :X\times_ZY\to X\times Y$ making the following diagram commute:
\[\begin{tikzcd}
	{X\times_ZY} && \\
	& {X\times Y} & X \\
	& Y
	\arrow["\varphi"{description}, dashed, from=1-1, to=2-2]
	\arrow[curve={height=-12pt}, from=1-1, to=2-3]
	\arrow[curve={height=12pt}, from=1-1, to=3-2]
	\arrow["{\pi_X}"', from=2-2, to=2-3]
	\arrow["{\pi_Y}", from=2-2, to=3-2]
\end{tikzcd}\]
    An element $a\in X\times_ZY$ is mapped to $\varphi(a)=(x,y)\in X\times Y$, and commutativity gives us that $x=p_X(a),y=p_Y(a)$. So by definition of $p_X,p_Y$, we have that whenever $\varphi(a)=(x,y)$, 
    \[
    f(x)=f(p_X(a))=g(p_Y(a))=g(y)
    \]
    We can therefore identify $X\times_ZY$ with the subset of $X\times Y$ which makes this equality true; that is,
    \[
    X\times_ZY=\{(x,y)\mid f(x)=g(y)\},
    \]
    where $p_X,p_Y$ are given by the usual projections restricted to this set.
\end{example}
\begin{example}[Specific situations in sets] \text{  }
    \begin{itemize}
        \item When $Z=\set{*}$ is a point, the only functions $f:X\to Z,g:Y\to Z$ are the trivial ones:
        \[\begin{tikzcd}
	       & X \\
	       Y & \bullet
	       \arrow[from=1-2, to=2-2]
	       \arrow[from=2-1, to=2-2]
        \end{tikzcd}\]
        So, using our characterisation above,
        \[
        X\times_{\set{*}}Y=\set{(x,y) \mid f(x)=g(y)}=X\times Y.
        \]
        That is, the fibered product over a point is the regular product.
        \item Let $Y\subseteq Z$ so that $i:Y\to Z$ is the inclusion map, and $f:X\to Z$ be any function. Using our characterisation, we have
        \[\begin{tikzcd}
        {\set{(x,y) \mid f(x)=y}} & X \\
	       Y & Z
	       \arrow[from=1-1, to=1-2]
	       \arrow[from=1-1, to=2-1]
	       \arrow["f", from=1-2, to=2-2]
	       \arrow[hook, from=2-1, to=2-2]
        \end{tikzcd}\]
        This set can be naturally identified with $f^{-1}(Y)$, the preimage of $Y$ under $f$. 
    \end{itemize}
\end{example}
\begin{remark}
    As for products, we should not expect the fibered product in the category of schemes to be the fibered product of the underlying sets + additional structure. Instead, what we should hope for is that fibered products respect $k$-points for fields (or rings) $k$, which aligns with the more canonical notion of points in algebraic geometry. \\
    Indeed, the universal property of the fibered product says that for any scheme $T$, 
    \[
    \mathrm{Maps}(T,X\times_ZY)=\mathrm{Maps}(T,X)\times_{\mathrm{Maps(T,Z)}}\mathrm{Maps(T,Y)}
    \]
    as sets. Denote the set of $k$-points of $X$ by $\mathrm{Maps}(\Spec k,X)\colonequals X(k)$. Then if $T=\Spec k$, this says that the $k$ points of $X\times_ZY$ correspond exactly to the fibered product $X(k)\times_{Z(k)}Y(k)$ (as sets). 
\end{remark}
All throughout the discussion on fibered products so far, it has not been clear that they actually exist! It would not be very fruitful to talk about them if they didn't, so:
\begin{theorem}
    Fibered products of schemes exist.
\end{theorem}
\begin{proof}
    We prove this when $X,Y,Z$ are affine schemes $\Spec A,\Spec B,\Spec R$, with maps $\Spec A\to \Spec R,\Spec B\to \Spec R$. We want $\square$ which makes the following diagram commute:
\[\begin{tikzcd}
	\bigcirc && \\
	& \square & {\Spec A} \\
	& {\Spec B} & {\Spec R}
	\arrow[dashed, from=1-1, to=2-2]
	\arrow[curve={height=-12pt}, from=1-1, to=2-3]
	\arrow[curve={height=12pt}, from=1-1, to=3-2]
	\arrow[from=2-2, to=2-3]
	\arrow[from=2-2, to=3-2]
	\arrow[from=2-3, to=3-3]
	\arrow[from=3-2, to=3-3]
\end{tikzcd}\]
    Since morphisms of affine schemes are exactly morphisms of their global sections in the other direction, this amounts to finding a ring $\square$ which satisfies the following universal property for any ring $\bigcirc$:
\[\begin{tikzcd}
	\bigcirc && \\
	& \square & A \\
	& B & R
	\arrow[dashed, from=2-2, to=1-1]
	\arrow[curve={height=12pt}, from=2-3, to=1-1]
	\arrow[from=2-3, to=2-2]
	\arrow[curve={height=-12pt}, from=3-2, to=1-1]
	\arrow[from=3-2, to=2-2]
	\arrow[from=3-3, to=3-2]
	\arrow[from=3-3, to=2-3]
\end{tikzcd}\]
    Such a $\square$ exists, and we will denote it by $A\otimes_R B$ --- the tensor product of $A$ and $B$ over $R$. To define it, consider (as we did for the product)
    \[
    G\colonequals \set{a\otimes b \mid a\in A,b\in B}.
    \]
    Let $A\otimes_RB$ be the free $R$-module generated by $G$, modulo the following relations:
    \[
    \begin{split}
        a_1 \otimes b + a_2\otimes b - (a_1 +a_2)\otimes b,&\qquad  \forall a_1, a_2 \in A,\quad  b \in B,\\
        a\otimes b_1 + a\otimes b_2 - a \otimes (b_1 + b_2), &\qquad \forall a \in A, \quad  b_1,b_2 \in B, \\
        (ra)\otimes b-r(a\otimes b), &\qquad \forall r\in R, \quad a\in A, \quad b\in B, \\
        a\otimes (rb)-r(a\otimes b), &\qquad \forall r\in R, \quad a\in A, \quad b\in B.
    \end{split}
    \]
    Here multiplication of an element in $A$ or $B$ by $r\in R$ is taken to mean multiplication by the image of $r$ under the maps $R\to A$ or $R\to B$ respectively. \\
    As we did before, we give $A\otimes_RB$ a ring structure by letting
    \[
    (a_1\otimes b_1)\cdot (a_2\otimes b_2)\colonequals a_1a_2\otimes b_1b_2,
    \]
    and extending linearly (the fact that this gives a well-defined ring structure needs checking!). Then the maps $A\to A\otimes_RB, B\to A\otimes_RB$ are given by $a\mapsto a\otimes1,b\mapsto 1\otimes b$ respectively. One can also check easily that these are ring homomorphisms. \\
    Moreover, we have from the given relations that for $r\in R$,
    \[
    r\otimes 1=r(1\otimes 1)=1\otimes r,
    \]
    and so the maps $R\to A\to A\otimes_RB$ and $R\to B\to A\otimes _RB$ are equal. That is, the following diagram commutes:
\[\begin{tikzcd}
	{A\otimes_RB} & A \\
	B & R
	\arrow[from=1-2, to=1-1]
	\arrow[from=2-1, to=1-1]
	\arrow[from=2-2, to=1-2]
	\arrow[from=2-2, to=2-1]
\end{tikzcd}\]
    Finally, let $f:A\to T$ and $g:B\to T$ be maps such that the compositions $R\to A\to T$ and $R\to B\to T$ are equal. Then there is a unique map $\varphi:A\otimes_RB\to T$ such that the composition $A\to A\otimes_RB\to T$ is $f$ and $B\to A\otimes_RB\to S$ is $g$. This is given by
    \[
    \varphi(a\otimes b)\colonequals f(a)\cdot g(b),
    \]
    and extending linearly. Once again, it is easy to check that this is a well-defined homomorphism, and that it is the unique one with this property. \\
    As a result, we can conclude that $\Spec A\times_{\Spec R}\Spec B$ exists, and is equal to $\Spec(A\otimes_RB)$. 
\end{proof}
Let's look at some examples of tensor products and how to compute them. Throughout, whenever we take the tensor product of $A,B$ over $R$, we assume we have maps $R\to A,B$.
\begin{example}\text{  }
    \begin{itemize}
        \item We claim that
        \[
        R\otimes_RS=S,
        \]
        where the map right to left is given by $r\otimes s\mapsto s$, and its inverse is $1\otimes s\mapsfrom s$. It's easy to check these are ring homomorphisms. \\
        Obviously these are not literally equal (even as sets!), but there is a natural canonical way to identify them. In general, it is often simpler to think of two objects related in this way as being equal, and many mathematicians choose to do so. 
        \item Considering polynomial rings, we have
        \[
        R[x]\otimes_R S=S[x],
        \]
        via the maps $f(x)\otimes s\mapsto sf(x)$ and $1\otimes s\mapsfrom s, x\otimes 1\mapsfrom x$ (recall that a map out of $S[x]$ is exactly a map out of $S$ together with an assignment of $x$).
    \end{itemize}
\end{example}
\begin{exercise}
    Find the maps that realise the following natural identifications:
    \begin{enumerate}
        \item For $I\lhd R$ and ideal, we have
        \[
        R/I\otimes_RS=S/IS,
        \]
        where $IS$ is the ideal of $S$ generated by the image of $R$.
        \item For $f\in R$, we have
        \[
        R_f\otimes_RS=S_f,
        \]
        where on the right we localise at the image of $f$ in $S$. 
        \item For $f_1,\ldots,f_m\in R[x]$, we have
        \[
        R[x_1,\ldots,x_n]/(f_1,\ldots,f_m)\otimes_RS=S[x_1,\ldots,x_n]/(f_1,\ldots,f_m),
        \]
        where on the right we take $f_i$ to be polynomials in $S[x_1,\ldots,x_n]$.
    \end{enumerate}
\end{exercise}
\begin{remark}
    These results all display a shortcut for how to think about the tensor product. When $A$ is obtained from $R$ by a particular algebraic construction, the tensor product $A\otimes_RB$ often corresponds to performing the same construction over $B$. Informally, tensoring with $B$ replaces the coefficients in $R$ with coefficients in $B$, while preserving the structure of the construction.
\end{remark}
\begin{example}
    \[
    \Z[x,y]/(x^2-2y)\otimes_\Z\Q=\Q[x,y]/(x^2-2y)
    \]
\end{example}
\begin{corollary}
    If $A$ is a finitely generated $R$-algebra, then $A\otimes_RB$ is also a finitely generated $R$-algebra. Moreover, if $A,B,R$ are all finitely generated $k$-algebras, then $A\otimes_RB$ is a finitely generated $k$-algebra.
\end{corollary}
\begin{example}
    Let's now consider the fiber product of the following configuration:
\[\begin{tikzcd}
	& {\Spec\C[x]} \\
	{\Spec\C} & {\Spec\C[y]}
	\arrow["{0\mapsfrom y}", from=2-1, to=2-2]
	\arrow["{y\mapsto x^2}"', from=1-2, to=2-2]
\end{tikzcd}\]
    where the maps are defined by the corresponding maps of global sections in the other direction. Everything here is affine, so this corresponds to the following diagram of rings:
\[\begin{tikzcd}
	& {\C[x]} \\
	\C & {\C[y]}
	\arrow["{y\mapsto x^2}"', from=2-2, to=1-2]
	\arrow["{0\mapsfrom y}"', from=2-2, to=2-1]
\end{tikzcd}\]
    To find the fibered product of the original setup, we need to compute the tensor product $\C\otimes_{\C[y]}\C[x]$. Writing $\C=\C[y]/(y)$ (which is how $\C$ is attained from $\C[y]$ in our situation), we can use our shortcuts to yield
    \[
    \C\otimes_{\C[y]}\C[x]=\C[y]/(y)\otimes_{\C[y]}\C[x]=\C[x]/(y)\C[x]
    \]
    The image of $y$ in $\C[x]$ is $(x^2)$, and so
    \[
    \C\otimes_{\C[y]}\C[x]=\C[x]/(y)\C[x]=\C[x]/(x^2).
    \]
    Thus the fibered product is $\Spec\C[x]/(x^2)$. Note that we could have viewed $\C[x]$ as being attained from $\C[y]$ via the construction $\C[x]=\C[y][x]/(y-x^2)$, in which case we would have
    \[
    \C\otimes_{\C[y]}\C[x]=\C\otimes_{\C[y]}\C[y][x]/(y-x^2)=\C[x]/(0-x^2)=\C[x]/(x^2),
    \]
    leading to the same result. \\
    If we instead replace the bottom map $y\mapsto 0$ by $y\mapsto 1$, then we have the following computation:
    \[
    \C\otimes_{\C[y]}\C[x]=\C[y]/(y-1)\otimes_{\C[y]}\C[x]=\C[x]/(y-1)\C[x]=\C[x]/(x^2-1),
    \]
    so the fibered product is $\Spec\C[x]/(x^2-1)$.
\end{example}
\begin{definition}[Scheme theoretic fiber]
    Let $k$ be a field, $f:X\to Y$ a morphism of schemes, and $\Spec k\to Y$ a $k$-point of $Y$. The \textit{scheme theoretic fiber} is defined to be the fibered product $X\times_Y\Spec  k$:
\[\begin{tikzcd}
	{X\times_Y\Spec k} & X \\
	{\Spec k} & Y
	\arrow[from=1-1, to=1-2]
	\arrow[from=1-1, to=2-1]
	\arrow["f", from=1-2, to=2-2]
	\arrow[from=2-1, to=2-2]
\end{tikzcd}\]
\end{definition}
In the previous example we computed the scheme theoretic fibers of the map $\Spec \C[x]\to \Spec \C[y]$ induced by $y\mapsto x^2$ at the $\C$-points $0$ and $1$. What we found was that the fiber at $0$ saw a sort of squashed infinitesimal fiber corresponding to $\Spec\C[x]/(x^2)$, which has a single $\C$-point which we like to think of as possessing some `fuzz' around it. On the other hand, the fiber at $1$ saw $\pm1$, corresponding to $\Spec\C[x]/(x^2-1)$ which has two $\C$-points. This is somewhat reminiscent of a double-sheeted covering of the complex plane, where a point has two distinct preimages, but at the branch point $0$ the sheets come together, producing a single point with that kind of `fuzz'.
\end{document}
